\section{The junction-vertex quotient, ultra hubs, and the induction}\label{s7:sec}

This section supplies the J-consumer of Definition~\ref{s6:defJconsumer}, bounds
the size of the quotient graphs it produces, and closes the argument by the
layered quotient induction. The mathematics is a \candidate{} argument: it has
been checked only by AI referees and has not been verified by human experts (see
Remark~\ref{s1:remStatus} and the final subsection of this section).

\paragraph{Setting of this section.}
Throughout, $\Dstar$ satisfies \Gc{1}--\Gc{4}, $G$ is a graph with $n\ge\Nzero$
vertices and $d_1\ge\Dstar$, a valid $\HBtp$ run on $G$ is fixed
(Proposition~\ref{s2:propExists}), a designation $\delta$ is fixed
(Definition~\ref{s6:defDesign}), and there are no VX-parts
($\mathcal{V}=\emptyset$), so that every standalone pre-part is a MIX-part. The
rounds of the run are $1,\dots,R$; the round step below acts at the rounds
$3\le l\le R$ (rounds $1$ and $2$ have no classed ports). For a classed port $u$
of round $l$ (that is, $u\in Q_Z$ for some $Z\in\Std_l$) we write $Y(u)$ for its
class and $r(u)$ for the round of $Y(u)$; since $Y(u)\in\anc_l(u)$ we have
$r(u)\le l-2$ and $u\in V(Y(u))$. For an ancestor $Y$ of round $r$, $H_Y$ is its
expander ($X_Y$ if $Y$ is light, $X^0_Y$ if $Y$ is standalone), $s_Y$ is $s_r/2$
or $s_r$ respectively, and $p_Y=1/(2\klend(Y))$ (Definition~\ref{s3:defCOL}).
The functions $\epsOne$ and $\epsTwo$ defined below contain the function
$\epsA=31\eps/(\Cp\log\log\Dstar)$ of Lemma~\ref{s2:lemTower}(e), which depends
on $\Dstar$ only and satisfies $\sum_{l\le R}15.2\eps/\log P_l\le\epsA$; the
function $\epsX$ does not contain $\epsA$.

%% ----------------------------------------------------------------------
\subsection{Round-split pools and candidate junctions}

\begin{definition}[Round-split pool labels]\label{s7:defPool}
Every vertex $v$ of $G$ independently draws a \emph{pool label}
\[
\plab(v)\in\{\bot\}\cup\{(l,r):3\le l\le R,\ 1\le r\le l-2\},\qquad
\Prob\bigl(\plab(v)=(l,r)\bigr)=q_l\,\pi_{l,r},
\]
where $q_l\defeq M_l^{-2}$ and $\pi_{l,r}\defeq 2^{-(l-1-r)}$, and
$\Prob(\plab(v)=\bot)\defeq 1-\sum_{l,r}q_l\pi_{l,r}$. This is a probability
distribution: for each $l$ we have $\sum_{r=1}^{l-2}\pi_{l,r}=1-2^{-(l-2)}<1$,
and $\sum_l q_l\le\sum_l M_l^{-1}\le 2/\Dstar<1$ by
Lemma~\ref{s2:lemTower}(b). Put
\[
\Pool_{l,r}\defeq\plab^{-1}(l,r),\qquad
\Pool_l\defeq\bigcup_{r=1}^{l-2}\Pool_{l,r}\ \ \text{(a disjoint union)},
\]
so that $\Prob(v\in\Pool_l)=q_l(1-2^{-(l-2)})\le q_l$ for every $v$. The sets
$\Pool_l$ ($3\le l\le R$) are pairwise disjoint, because every vertex has one
label. For $w\in\Pool_l$ let $r(w)$ denote the unique $r$ with
$w\in\Pool_{l,r}$. The pool labels are stage (1d) of the randomness schedule.
They are independent of all other stage-1 data, in particular of the colourings
and labels of Definition~\ref{s3:defCOL}.
\deps{\ref{s2:lemTower}, \ref{s3:defCOL}}
\end{definition}

\begin{definition}[Candidates and JV-bad ports]\label{s7:defCand}
Let $u$ be a classed port of round $l\ge3$, with class $Y\defeq Y(u)$ of round
$r\defeq r(u)\le l-2$. Its \emph{candidate set} is
\[
\Cand_l(u)\defeq\{\,w\in\Pool_{l,r(u)}:\ uw\in\LJV_{Y,l}\,\}.
\]
Since $\LJV_{Y,l}\subseteq\Lend_Y\subseteq E(H_Y)$ and $H_Y$ is a graph on
$V(Y)$, automatically $\Cand_l(u)\subseteq N_{H_Y}(u)$ and
$N_{H_Y}(u)\subseteq V(Y)$. The
candidate sets are always taken in the class graph $H_{Y(u)}$ itself; this is
used in Lemma~MULT below. Put
\[
\Hcd_l\defeq\frac{\lambda_{l-2}^{95}}{8M_l^{2}}.
\]
The port $u$ is \emph{JV-bad} if
\[
|\Cand_l(u)|<\tfrac12\,\Ex|\Cand_l(u)|,
\]
and \emph{JV-good} otherwise. Here $\Ex$ is the unconditional expectation over
stage~1, which is a number determined by the run and $\delta$. The definition
applies to every classed port, whether or not it lies in $\Lost_l$. JV-badness
is a deterministic function of stage~1; it is one of the stage-2 statuses. The
number $\Hcd_l$ is not a threshold of the definition: it is a lower bound for
the threshold, since $\tfrac12\Ex|\Cand_l(u)|\ge\Hcd_l$ for every classed port
$u$ of round $l$ (Lemma~\ref{s7:lemCand}(ii) below). Consequently every JV-good
port has $|\Cand_l(u)|\ge\Hcd_l$ (Lemma~\ref{s7:lemCand}(iv)).
\deps{\ref{s7:defPool}, \ref{s3:defCOL}, \ref{s6:defDesign}}
\fixes{\RTa{M4}}
\end{definition}

\begin{lemma}[Candidate counts]\label{s7:lemCand}
Let $u$ be a classed port of round $l\ge3$ with class $Y$ of round $r$. Then:
\begin{enumerate}
\item[(i)] $|\Cand_l(u)|$ is a sum of independent Bernoulli variables, one for
each $w\in N_{H_Y}(u)$, each with success probability $q_l\pi_{l,r}p_Y$;
\item[(ii)] its mean satisfies
\[
\Ex|\Cand_l(u)|=q_l\,\pi_{l,r}\,p_Y\,\deg_{H_Y}(u)
\;\ge\;\frac{\lambda_r^{96}\,2^{-(l-r)}}{M_l^{2}}
\;\ge\;2\Hcd_l ;
\]
\item[(iii)] $\Prob(u\text{ is JV-bad})\le\exp(-\Hcd_l/4)$;
\item[(iv)] if $u$ is JV-good then $|\Cand_l(u)|\ge\Hcd_l\ge2^{10}M_l^{10}$.
\end{enumerate}
\deps{\ref{s7:defCand}, \ref{s7:defPool}, \ref{s3:defCOL}, \ref{s3:lemCOLJV}, \ref{s2:propStructure}, \ref{s2:lemTower}, \ref{s1:citChernoffGen}}
\fixes{\RTa{M8}, \RTb{J7}}
\end{lemma}

\begin{proof}
(i) By definition,
\[
|\Cand_l(u)|=\sum_{w\in N_{H_Y}(u)}\ind{\plab(w)=(l,r)}\cdot\ind{uw\in\LJV_{Y,l}} .
\]
The pool labels are independent across vertices. The colours of the edges of
$H_Y$ are independent across edges (Definition~\ref{s3:defCOL}), and the edges
$uw$ for distinct $w$ are distinct. The two families are independent of each
other (Definition~\ref{s7:defPool}). So the summands for distinct $w$ are
functions of disjoint sets of independent variables, hence independent. For a
fixed $w$ the two indicators are independent, with
$\Prob(\plab(w)=(l,r))=q_l\pi_{l,r}$. Moreover $\Prob(uw\in\LJV_{Y,l})=p_Y$: the
edge $uw\in E(H_Y)$ lies in $\Lend_Y$ with probability $1/2$, and it then
receives the colour $l\in\IJV$ with probability $1/\klend(Y)$; note that
$r+2\le l\le R$, so $l$ is an index of $\IJV$.

(ii) The first equality follows from (i). Every $H_Y$ has minimum degree larger
than $s_Y\ge s_r/2$, and $s_r\ge\lambda_r^{100}$
(Proposition~\ref{s2:propStructure}(i)). Also $p_Y\ge\lambda_r^{-4}$
(Lemma~\ref{s3:lemCOLJV}). Hence
\[
\Ex|\Cand_l(u)|\ \ge\ M_l^{-2}\,2^{-(l-1-r)}\,\lambda_r^{-4}\,\frac{\lambda_r^{100}}{2}
\;=\;\frac{\lambda_r^{96}\,2^{-(l-r)}}{M_l^{2}} .
\]
Since $2\Hcd_l=\lambda_{l-2}^{95}/(4M_l^2)$, the second inequality is equivalent
to
\begin{equation}\label{s7:eqCandExp}
\lambda_r^{96}\ \ge\ 2^{\,l-r-2}\,\lambda_{l-2}^{95}.
\end{equation}
By Lemma~\ref{s2:lemTower}(a), $\lambda_r\ge\lambda_{l-2}$ because $r\le l-2$.
By Lemma~\ref{s2:lemTower}(d),
$\lambda_r\ge2^{2\logs d_r+2}\ge2^{R-r}\ge2^{l-r}\ge2^{l-r-2}$. Multiplying
$\lambda_r^{95}\ge\lambda_{l-2}^{95}$ by $\lambda_r\ge2^{l-r-2}$ gives
\eqref{s7:eqCandExp}. (The chain of exponents is
$2\logs d_r+2\ge R-r\ge l-r$; only the exponent $l-r-2$ is needed.)

(iii) Put $\mu\defeq\Ex|\Cand_l(u)|$. By definition, $u$ is JV-bad exactly
when $|\Cand_l(u)|<\mu/2$. By (i) and the lower-tail Chernoff
bound for sums of independent indicators, with deviation parameter $1/2$
(Cited result~\ref{s1:citChernoffGen}),
$\Prob(|\Cand_l(u)|<\mu/2)\le\exp(-\mu/8)\le\exp(-\Hcd_l/4)$, where the last
step uses $\mu/2\ge\Hcd_l$ from (ii).

(iv) If $u$ is JV-good then $|\Cand_l(u)|\ge\tfrac12\Ex|\Cand_l(u)|\ge\Hcd_l$,
by the definition of JV-good and (ii); this is the first inequality. The second
is the line
``$\Hcd_l\ge2^{10}M_l^{10}$'' of the table of Lemma~\ref{s3:lemCOLJV}. It also
follows directly from $M_l\le\lambda_{l-2}^{1.6}$
(Lemma~\ref{s2:lemTower}(b)): then
$8M_l^2\cdot2^{10}M_l^{10}=2^{13}M_l^{12}\le2^{13}\lambda_{l-2}^{19.2}\le\lambda_{l-2}^{95}$,
because $\lambda_{l-2}\ge M_l^{1/1.6}\ge2^{25}$.
\end{proof}

%% ----------------------------------------------------------------------
\subsection{The randomness schedule and the round step}

\begin{definition}[The randomness schedule]\label{s7:defSchedule}
All random choices of the construction are made in the following order,
summarized in Table~\ref{s7:tabSchedule}.
\begin{itemize}
\item \emph{Stage 1} consists of four mutually independent families:
(1a) the COL-JV colourings of Definition~\ref{s3:defCOL}(i)--(iii);
(1b) the vertex choices and sublabels of Definition~\ref{s5:defZones};
(1c) the JS labels of Definition~\ref{s3:defCOL}(iv);
(1d) the pool labels of Definition~\ref{s7:defPool}.
\item \emph{Selection} (1e), which is not a random family: a stage-1 outcome is
then fixed in the event $\{\Xp\le3\,\Ex\Xp\}$. Here
$\Xp\ge0$ is the stage-1 functional introduced later in this section, a
deterministic function of the run, $\delta$ and stage~1. After (1e), stage~1 is
a fixed outcome: no later step of the construction uses the law of stage~1 or
conditions on the event of (1e); the numbers $\Ex\Xp$ and $\Ex|\Cand_l(u)|$
that appear later are unconditional expectations, fixed in advance.
\item \emph{Stage 2} is deterministic given stage~1: demotion, parent-bad and
$\Bad$ (Definition~\ref{s5:defStages}); lend-bad, $\Lost$, $\Ret$ and $\Qs$
(Definition~\ref{s6:defLending}); JV-bad (Definition~\ref{s7:defCand}).
\item \emph{Stage 3} is drawn part by part: the TPV labels of every standalone
pre-part, the P-vortex labels of every non-demoted light part, and the VX$^+$
labels of every demoted light part. Given stage~1, the label families of distinct
parts are independent. Each part's labels are fixed inside that part's good event.
\item \emph{Rounds} $l=R,R-1,\dots,3$. The round randomness $\xi_l$ consists of
the following mutually independent variables:
a uniformly random bijection $\eta_h:[4M_l]\to[4M_l]$ for every vertex $h$;
a uniformly random $3$-subset $\zeta_{h,u}\subseteq[4M_l]$ for every ordered pair
$(h,u)$ of distinct vertices;
and a uniformly random linear order $\prec_u$ of $V(G)$ for every vertex $u$.
The dependence of $\eta_h,\zeta_{h,u},\prec_u$ on $l$ is suppressed. The draw
$\xi_l$ is fresh: it is independent of stage~1, of stage~3 and of all $\xi_{l'}$
with $l'>l$.
\end{itemize}
The \emph{past} $\Past_l$ at round $l$ is the fixed outcome of stage~1, of
stage~3 and of the draws $\xi_{l'}$ ($l'>l$), together with everything these
determine: in particular the decompositions fixed at the rounds $l'>l$ and the
set $J_l$. It is a fixed outcome, not a random $\sigma$-algebra. We write
$\Ex[\,\cdot\mid\Past_l]$ and $\Prob(\,\cdot\mid\Past_l)$ for expectation and
probability over $\xi_l$ with $\Past_l$ fixed. At each round, $\xi_l$ is chosen in
the event of step (f) of the round step. All draws in stage~1, stage~3 and the
rounds are mutually independent. The outcomes actually used are chosen inside
events (stage~1 in $\{\Xp\le3\,\Ex\Xp\}$, stage~3 in the good events, each
$\xi_l$ in the event of step (f)); after each choice only deterministic
properties of the chosen outcome are used, and every later probability is over
fresh variables with the past fixed.
\deps{\ref{s3:defCOL}, \ref{s5:defZones}, \ref{s7:defPool}, \ref{s6:defLending}, \ref{s5:defStages}, \ref{s7:defCand}, \ref{s4:lemTPV}, \ref{s4:lemPV}, \ref{s4:thmVXp}, \ref{s5:lemChild}, \ref{s5:lemDemoted}}
\fixes{\RTb{J5}, \RTa{M4}}
\end{definition}

\begin{table}[ht]
\centering
\small
\begin{tabular}{|p{1.5cm}|p{5.9cm}|p{3.6cm}|p{3.8cm}|}
\hline
\textbf{Stage} & \textbf{Drawn or decided} & \textbf{Depends on} & \textbf{How it is fixed}\\
\hline
1a & COL-JV colourings of every ancestor & independent per ancestor and per edge & part of the stage-1 outcome\\
1b & vertex choices and sublabels (zones) & independent per vertex & part of the stage-1 outcome\\
1c & JS labels $\lab_{Y,l}(y)$ & independent & part of the stage-1 outcome\\
1d & pool labels $\plab(v)$ & independent per vertex & part of the stage-1 outcome\\
1e & choice of the stage-1 outcome & the functional $\Xp$ & $\Xp\le3\,\Ex\Xp$ (probability $\ge2/3$)\\
\hline
2 & demotion, parent-bad, lend-bad, $\Lost$, $\Ret$, $\Qs$, JV-bad & deterministic in stage~1 & determined\\
\hline
3 & TPV, P-vortex and VX$^+$ labels & independent per part & inside each part's good event (probability $\ge1/2-o(1)$)\\
\hline
$l=R,\dots,3$ & $\xi_l=(\eta_h,\zeta_{h,u},\prec_u)$: HUB lists and injections & fresh draw, $\Past_l$ fixed & event of step (f) (probability $\ge1/2$)\\
\hline
\end{tabular}
\caption{The randomness schedule (Definition~\ref{s7:defSchedule}).}
\label{s7:tabSchedule}
\end{table}

\begin{construction}[The JV$^{+*}$ round step at round $l\ge3$; the J-consumer]\label{s7:consRound}
Fix $3\le l\le R$ and the past $\Past_l$. The input is the set
$J_l=\Jlost_l\cup\Jhub_l\cup\Jfr_l\cup\Jpar_l$ produced by JS-LC at round $l$,
with the properties \Jp{1}--\Jp{3} and (i)--(v) of Lemma~\ref{s6:lemJplus}.

\emph{Fixed rules.} The rules and orders called fixed in steps (b)--(e), (g)
and (h) below are deterministic functions, chosen once and for all before any
randomness is drawn, with the following arguments.
\begin{itemize}
\item \emph{Functions of $\Past_l$ alone:} the pairing of the live $\Jfr$ items
at a fresh centre and the choice of its unpaired fresh leg, and the order in
which the PAR objects are coloured, in (b); the split of the live items of a
non-ultra hub into groups, together with the numbering of the groups, in (c).
In particular these rules do not read $\xi_l$.
\item \emph{Functions of $\Past_l$ and of the lists}, that is, of the variables
$\eta_h$ and $\zeta_{h,u}$ of $\xi_l$: the choice of the system of distinct
representatives in (d); the processing order and the order of $V(G)$ in (e1);
and the order $e_1,\dots,e_q$ of $E'(u)$ in (e2).
\item None of the rules of steps (b)--(e) reads the orders $\prec_u$ of
$\xi_l$; within steps (a)--(f), the orders $\prec_u$ of $\xi_l$ enter only
through the assignment $e_i\mapsto w_i$ of (e2) (and through what is computed
from it). Through $\Past_l$ the rules may depend on the draws
$\xi_{l'}$ of the rounds $l'>l$, orders included, and in general they do,
since $J_l$ does.
\item The rank order in (g) and the rule choosing $\Dec_l$ in (h) are not
restricted: they are arbitrary deterministic functions of $\Past_l$ and
$\xi_l$. They necessarily depend on $\xi_l$, since the layers and $Q_l$ do.
\end{itemize}
Rules with these arguments exist, for instance lexicographically first choices
for fixed orders of the finite sets involved (the greedy colouring of (b)
succeeds for every order, by part (i) of the next lemma, and a split into groups
as in (c) always exists). Every statement of this section about the
round step holds for every choice of the fixed rules that obeys these
restrictions. The restrictions are used in the claims of steps (c) and (e2), in
Lemma~\ref{s7:lemWellDef}(iii) (the grouping does not read $\eta_h$), and in
Lemmas~\ref{s7:lemCC}(i), \ref{s7:lemUltra}(i) and \ref{s7:lemPay}(c) (nothing
before the assignment $e_i\mapsto w_i$ reads the orders $\prec_u$), as well as
in the opening paragraph of the subsection ``Quotient size and payments'' and
in the proof of Lemma~\ref{s7:lemUHsplit}(ii) (the copies of non-ultra hubs are
deterministic given the lists).

\emph{Items.} The \emph{items} are the edges of $J_l\setminus\Jlost_l$. The
\emph{vertices} of an item are its port end or ends, which lie in $\Qs_Z$, and
its centre. The centre is a hub $h\in D_l$ for a $\Jhub$ item and a fresh centre
$x\in F_Z$ for a $\Jfr$ item. A $\Jpar$ item $u$--$v$ has two port ends and no
centre. A $\Jhub$ item is written $(h,u)$, with $h$ its hub and $u$ its port.

\smallskip\noindent
(a) \emph{Payments bounded by stage-1 weights.} The following are \emph{paid},
that is, output as single edges:
(a1) every edge of $\Jlost_l$;
(a2) every item with a vertex in $\Pool_l$;
(a3) every item with a JV-bad port end.
The remaining items are \emph{live}. A vertex is \emph{live} if it is a vertex of
a live item, and $\Live$ denotes the set of live vertices. By (a2),
$\Live\cap\Pool_l=\emptyset$. The sets paid in (a1)--(a3) depend on $J_l$, hence
on the past; only their sizes are bounded by stage-1 functionals
(this is proved in the payment lemma below).

\smallskip\noindent
(b) \emph{PAR objects.} Every live $\Jpar$ item $u$--$v$ is a \emph{PAR object}
with \emph{ends} at $u$ and at $v$; here $u,v\in\Qs_Z$ and $Y(u)\ne Y(v)$. At each
fresh centre $x$, the live $\Jfr$ items at $x$ are paired, by a fixed rule, into
\emph{cherries} $u$--$x$--$u'$. Here $u\ne u'$ because $G$ is simple, so
distinct items at $x$ have distinct port ends. A cherry is a PAR object with ends
at $u$ and $u'$ and \emph{middle} $x$. If the number of live $\Jfr$ items at $x$
is odd, one of them, chosen by a fixed rule, is paid: it is the \emph{unpaired
fresh leg} of $x$. The PAR objects are then coloured greedily, in a fixed order,
with colours from $[3M_l]$, so that two PAR objects sharing a port, or sharing a
middle, receive different colours.

\smallskip\noindent
(c) \emph{HUB lists.} For a hub $h\in D_l$ let $\clive_h$ be its number of live
$\Jhub$ items. Put
\[
\thult_l\defeq\bigl\lfloor M_l\Hcd_l/7\bigr\rfloor,\qquad \KHUB_l\defeq4M_l .
\]
The hub $h$ is \emph{ultra} if $\clive_h>\thult_l$, and \emph{non-ultra}
otherwise. For a non-ultra $h$ with $\clive_h\ge1$, its live items are split, by
a fixed rule, into
\[
k_h\defeq\bigl\lceil 8\clive_h/\Hcd_l\bigr\rceil
\]
\emph{groups} of at most $\lceil\Hcd_l/8\rceil$ items each. This is possible
because $k_h\lceil\Hcd_l/8\rceil\ge\clive_h$. The $i$-th group receives the
\emph{list} $\{\eta_h(3i-2),\eta_h(3i-1),\eta_h(3i)\}\subseteq[4M_l]$; this needs
$3k_h\le4M_l$, which is part (ii) of the next lemma. The lists of the
groups of $h$ are pairwise disjoint, and, since the groups and their numbering
are functions of $\Past_l$ while $\eta_h$ is independent of $\Past_l$, they form
a uniformly random sequence of
$k_h$ pairwise disjoint $3$-subsets of $[4M_l]$. Lists of distinct hubs are
independent. For an ultra $h$, every live item $(h,u)$ is its own group, with
list $\zeta_{h,u}$.

\smallskip\noindent
(d) \emph{SDR.} Every port $u$ carrying at least one live hub item chooses, by a
fixed rule, pairwise distinct colours for its live hub items, each item's colour
taken from the list of its group (a system of distinct representatives). If no
such choice exists, all live hub items of $u$ are paid (an \emph{SDR failure} at
$u$). The items that receive a colour are the \emph{coloured hub items}. In
particular, every port carries at most one coloured hub item of each HUB colour.

\smallskip\noindent
(e) \emph{Junctions.} Initially $\Used(u)\defeq\emptyset$ for every port $u$ and
$\Used_\kappa(h)\defeq\emptyset$ for every hub $h$ and every $\kappa\in[4M_l]$.
\begin{itemize}
\item[(e1)] The coloured hub items of non-ultra hubs are processed in a fixed
order. The item $(h,u)$ of colour $\kappa$ receives as its \emph{junction} the
first vertex $w$, in a fixed order of $V(G)$, of
$\Cand_l(u)\setminus\Used(u)\setminus\Used_\kappa(h)$. Then $w$ is added to
$\Used(u)$ and to $\Used_\kappa(h)$.
\item[(e2)] For every port $u$ let $E'(u)$ be the set of the ends at $u$ of the
PAR objects and of the coloured hub items $(h,u)$ with $h$ ultra, listed in a
fixed order $e_1,\dots,e_q$. Let $w_1\prec_u w_2\prec_u\cdots$ be the elements of
$\Cand_l(u)\setminus\Used(u)$ in the order $\prec_u$; the end $e_i$ receives the
junction $w_i$. This is well defined by part (v) of the next lemma. With
$\Past_l$ and the lists fixed, the map $e_i\mapsto w_i$ is a uniformly random
injection of $E'(u)$ into $\Cand_l(u)\setminus\Used(u)$, and these injections are
independent over ports: by the restrictions on the fixed rules, $\Used(u)$, the
set $E'(u)$ and its order $e_1,\dots,e_q$ are determined by $\Past_l$ and the
lists, and the orders $\prec_u$ are independent of both. No payment is made for
coincidences of junctions at ultra hubs.
\item[(e3)] \emph{Loops.} A PAR object whose two ends received the same junction
is paid. A looped $\Jpar$ item costs one single edge; a looped cherry costs its
two edges.
\end{itemize}
Every end of a coloured hub item or of an unpaid PAR object now has a junction
$w$. It is a candidate of the end's port $u$, so $w\in\Cand_l(u)\subseteq\Pool_l$
and $uw\in\LJV_{Y(u),l}$. The edge $uw$ is the \emph{junction edge} of that end.

\smallskip\noindent
(f) \emph{Markov choice.} Let $\payrd_l$ be the number of edges paid in (d) and
(e3), and let $\copies_l\defeq|V(Q_l)|$, with $Q_l$ as in (g). Both are
functions of $\Past_l$ and $\xi_l$. The round randomness $\xi_l$ is chosen in the
event
\[
\bigl\{\copies_l\le4\,\Ex[\copies_l\mid\Past_l]\bigr\}\cap
\bigl\{\payrd_l\le4\,\Ex[\payrd_l\mid\Past_l]\bigr\}.
\]
By Markov's inequality in the form~(b) of Cited result~\ref{s1:citMarkov},
taken conditionally as in its item~(c) ($\xi_l$ is independent of the variables
whose outcome $\Past_l$ fixes, Definition~\ref{s7:defSchedule}), this event has
probability at least $1/2$ given $\Past_l$ (each of the two events fails with
probability at most $1/4$).

\smallskip\noindent
(g) \emph{Layers, sub-layers and the quotient.} For $\kappa\in[3M_l]$ the
\emph{PAR layer} $\BP_\kappa$ is the multigraph on the vertex set
$\{[w]:w\in\Pool_l\}$ of formal copies. It has one edge $[w_1][w_2]$ for every
unpaid PAR object of colour $\kappa$ whose ends have the junctions $w_1$ and
$w_2$; here $w_1\ne w_2$ by (e3). For $\kappa\in[4M_l]$ the \emph{HUB layer}
$\BH_\kappa$ is the bipartite multigraph with sides $D_l$ and
$\{[w]:w\in\Pool_l\}$. It has one edge $h[w]$ for every coloured hub item $(h,u)$
of colour $\kappa$ whose junction is $w$. \emph{Rank split:} in every layer, the
edges joining the same two vertices are listed in a fixed order, and the $i$-th
edge of such a list has \emph{rank} $i$. For a layer $B$ and $\iota\ge1$, the
\emph{sub-layer} $B^{(\iota)}$ is the spanning subgraph of $B$ formed by its edges
of rank $\iota$. The \emph{quotient} $Q_l$ is the vertex-disjoint union of all
sub-layers $B^{\mathrm{P},(\iota)}_\kappa$ ($\kappa\in[3M_l]$, $\iota\ge1$) and
$B^{\mathrm{H},(\iota)}_\kappa$ ($\kappa\in[4M_l]$, $\iota\ge1$), with the
isolated vertices of each sub-layer deleted. Vertices of distinct sub-layers are
distinct vertices of $Q_l$; hence every cycle of $Q_l$ lies in one sub-layer.

\smallskip\noindent
(h) \emph{Decompose and lift.} Let $\Dec_l$ be a decomposition of $E(Q_l)$ into
$\f(Q_l)$ objects, chosen by a fixed deterministic rule (for instance the
lexicographically first one for a fixed order of $E(Q_l)$). Each member of
$\Dec_l$ is \emph{lifted} to objects of $G$ as follows.
\begin{itemize}
\item A single edge of a PAR sub-layer is a PAR object. It is output as its one
($\Jpar$) or two (cherry) J-edges, each as a single edge. A single edge $h[w]$ of
a HUB sub-layer is a coloured hub item $(h,u)$, and it is output as the single
edge $hu$.
\item A cycle $[w_0][w_1]\cdots[w_{m-1}]$ of a PAR sub-layer, with $m\ge3$:
indices are taken mod $m$, and the edge $[w_i][w_{i+1}]$ is a PAR object $o_i$.
Let $p_i$ be the end of $o_i$ with junction $w_i$ and $q_i$ its end with junction
$w_{i+1}$; they are well defined because $w_i\ne w_{i+1}$. The lift is the closed
walk
\[
w_0\,p_0\,(x_0)\,q_0\,w_1\,p_1\,(x_1)\,q_1\,w_2\ \cdots\ q_{m-1}\,w_0 ,
\]
where $x_i$ is the middle of $o_i$ if $o_i$ is a cherry, and is omitted if $o_i$
is a $\Jpar$ item.
\item A cycle $h_0[w_0]h_1[w_1]\cdots h_{m-1}[w_{m-1}]h_0$ of a HUB sub-layer,
with $m\ge2$: the edge $h_i[w_i]$ is a coloured hub item $(h_i,u_i)$ with
junction $w_i$, and the edge $[w_i]h_{i+1}$ is a coloured hub item
$(h_{i+1},u'_i)$ with junction $w_i$ (indices mod $m$). The lift is the closed
walk
\[
h_0\,u_0\,w_0\,u'_0\,h_1\,u_1\,w_1\,u'_1\,h_2\ \cdots\ u'_{m-1}\,h_0 .
\]
\end{itemize}
Put $\Obj_l\defeq{}$(the single edges paid in (a), (b), (d) and (e3)) together
with (the lifted objects of all members of $\Dec_l$), and let $\LentJV_l$ be the
set of junction edges that lie on lifted cycles. A junction edge assigned in (e)
that does not lie on a lifted cycle is not used by the round step; it stays with
its owner as junk.
\deps{\ref{s6:lemJplus}, \ref{s6:defJconsumer}, \ref{s7:defCand}, \ref{s7:defPool}, \ref{s7:defSchedule}, \ref{s1:citMarkov}}
\fixes{\RTb{J5}, \RTb{J8}, \RTb{J4}, \RTa{M4}}
\end{construction}

\begin{lemma}[The round step is well defined]\label{s7:lemWellDef}
For every past $\Past_l$ ($3\le l\le R$) the following hold.
\begin{enumerate}
\item[(i)] \emph{PAR palette.} Every PAR object shares a port or a middle with
fewer than $3M_l$ other PAR objects, so the greedy colouring of (b) succeeds.
Consequently every port carries at most one PAR object of each colour, and
every fresh centre is the middle of at most one cherry of each colour.
\item[(ii)] \emph{Disjoint lists.} Every non-ultra hub $h$ has
$k_h\le\lceil8M_l/7\rceil$ and $3k_h\le4M_l$.
\item[(iii)] \emph{Cuckoo SDR.} For every port $u$,
$\Prob(\text{SDR failure at }u\mid\Past_l)\le(\KHUB_l)^{-4}=(4M_l)^{-4}$.
\item[(iv)] \emph{Greedy step.} Step (e1) always finds a junction. Two coloured
items of the same non-ultra hub $h$ with the same colour $\kappa$ receive
distinct junctions.
\item[(v)] \emph{Injections.} At every port $u$ carrying a live item,
$|\Cand_l(u)\setminus\Used(u)|\ge\Hcd_l-(M_l-1)\ge\Hcd_l/2\ge M_l>|E'(u)|$, so
step (e2) is well defined.
\end{enumerate}
\deps{\ref{s7:consRound}, \ref{s6:lemJplus}, \ref{s7:lemCand}, \ref{s3:lemCOLJV}, \ref{s1:citHall}, \ref{s2:defHBtp}}
\fixes{\RTb{I3}, \RTb{J3}, \RTb{J8}}
\end{lemma}

\begin{proof}
Throughout, $M_l\ge2^{40}$ by the definition of $M_l$, and
$\Hcd_l\ge2^{10}M_l^{10}$ by Lemma~\ref{s7:lemCand}(iv) (equivalently, by the
table of Lemma~\ref{s3:lemCOLJV}). By \Jp{2} of Lemma~\ref{s6:lemJplus}, every
port carries at most $M_l-1$ J-edges, hence at most $M_l-1$ live items. By
Lemma~\ref{s6:lemJplus}(ii), every fresh centre carries at most $M_l-1$ $\Jfr$
edges, hence is the middle of at most $(M_l-1)/2$ cherries.

(i) A PAR object has at most two ports and at most one middle. Each of its ports
carries at most $M_l-2$ further live items, and each further PAR object at that
port uses one of them. Its middle, if any, is the middle of at most
$(M_l-1)/2-1$ further cherries. So the object conflicts with at most
$2(M_l-2)+(M_l-1)/2<3M_l$ others. When it is coloured, fewer than $3M_l$ colours
are therefore blocked, and a colour of $[3M_l]$ remains. The consequences are
the defining property of the colouring.

(ii) If $h$ is non-ultra then $\clive_h\le\thult_l\le M_l\Hcd_l/7$, so
$k_h=\lceil8\clive_h/\Hcd_l\rceil\le\lceil8M_l/7\rceil$. Hence
$3k_h\le3(8M_l/7+1)=24M_l/7+3\le4M_l$, because $M_l\ge21/4$.

(iii) Fix $u$ and let $(h_1,u),\dots,(h_m,u)$ be its live hub items, so that
$m\le M_l-1\le K/4$ with $K\defeq\KHUB_l=4M_l$. These items are distinct edges at
$u$ and $G$ is simple, so the hubs $h_1,\dots,h_m$ are distinct. With $\Past_l$
fixed, it is determined which hubs are ultra, how the items are grouped, and
which group each item belongs to (the grouping rule of step (c) reads $\Past_l$
only, by the restrictions on the fixed rules in
Construction~\ref{s7:consRound}). The list of $(h_i,u)$ is a function of
$\eta_{h_i}$ if $h_i$ is non-ultra, and equals $\zeta_{h_i,u}$ if $h_i$ is ultra.
Since the $h_i$ are distinct, the $m$ lists are independent. Each list is a
uniformly random $3$-subset of $[K]$; for a non-ultra hub this holds because the
sequence of its lists is a uniformly random sequence of pairwise disjoint
$3$-subsets, so each single list is uniform. By Hall's theorem (Cited
result~\ref{s1:citHall}, applied with the $m$ items as the index set and, for
each item $(h_i,u)$, its list as the set of that item), an SDR fails to exist
only if some set of $s$ items has at most $s-1$ colours in the union of their
lists. Since every list has $3$ elements, this forces $s\ge4$. The union is
then contained in some $(s-1)$-subset of $[K]$. The union bound over $s$, over
the $s$-sets of items and over the $(s-1)$-subsets of $[K]$ gives
\[
\Prob(\text{SDR failure at }u\mid\Past_l)\ \le\ \sum_{s=4}^{m}T_s,\qquad
T_s\defeq\binom{m}{s}\binom{K}{s-1}\Bigl(\binom{s-1}{3}\Big/\binom{K}{3}\Bigr)^{s}.
\]
We use $\binom{m}{s}\le(em/s)^s\le(eK/(4s))^s$,
$\binom{K}{s-1}\le(eK/(s-1))^{s-1}$, and
\[
\binom{s-1}{3}\Big/\binom{K}{3}=\frac{(s-1)(s-2)(s-3)}{K(K-1)(K-2)}\le\Bigl(\frac{s-1}{K}\Bigr)^3 ,
\]
which holds because $(s-1-j)/(K-j)\le(s-1)/K$ for $j=1,2$ and $s\le K+1$.
Multiplying, and using $(s-1)^{2s+1}\le s^{2s+1}$,
\[
T_s\ \le\ e^{2s-1}\,4^{-s}\,s^{-s}\,(s-1)^{2s+1}\,K^{-s-1}\ \le\
a_s\defeq\frac1e\cdot\frac sK\cdot\Bigl(\frac{e^2s}{4K}\Bigr)^{s}.
\]
First, $a_4=4e^7K^{-5}<4400\,K^{-5}$. Second, for $s\ge4$,
\[
\frac{a_{s+1}}{a_s}=\frac{s+1}{s}\cdot\frac{e^2(s+1)}{4K}\cdot\Bigl(\frac{s+1}{s}\Bigr)^{s}
\le\frac54\cdot\frac{e^2(s+1)}{4K}\cdot e=\frac{5e^3}{16}\cdot\frac{s+1}{K},
\]
which is at most $\frac{5e^3}{16}(\frac1{13}+\frac1K)<0.49$ when $4\le s<K/13$ and
$K\ge10^4$. Hence $a_s\le2^{4-s}a_4$ for $4\le s\le\lceil K/13\rceil$, and the sum
of these terms is at most $2a_4<8800\,K^{-5}$. Third, for $K/13<s\le m\le K/4$ we
have $s/K\le1/4$, so $a_s\le\frac1{4e}(e^2/16)^s\le0.47^{s}$, and the sum of
these terms is at most $0.47^{K/13}/0.53\le2\cdot0.47^{K/13}$. Altogether
\[
\Prob(\text{SDR failure at }u\mid\Past_l)\ \le\ 8800\,K^{-5}+2\cdot0.47^{K/13}\ \le\ K^{-4}
\qquad(K\ge10^4),
\]
since $8800K^{-5}\le0.88K^{-4}$ and $2\cdot0.47^{K/13}\le0.1K^{-4}$ for
$K\ge10^4$. Here $K=4M_l\ge2^{42}$.

(iv) Items reach (e1) only if they are live, so their ports are JV-good and
$|\Cand_l(u)|\ge\Hcd_l$ (Lemma~\ref{s7:lemCand}(iv)). When the item $(h,u)$ of
colour $\kappa$ is processed, $\Used(u)$ contains one junction for each earlier
item at $u$, so $|\Used(u)|\le M_l-2$. The lists of the groups of $h$ are pairwise
disjoint, so all coloured items of $h$ with colour $\kappa$ lie in the one group
whose list contains $\kappa$. That group has at most $\lceil\Hcd_l/8\rceil$
items, so $|\Used_\kappa(h)|\le\lceil\Hcd_l/8\rceil-1<\Hcd_l/8$. Hence at least
$\Hcd_l-(M_l-2)-\Hcd_l/8>0$ choices remain. The second statement holds because
each junction chosen for $h$ in colour $\kappa$ is added to $\Used_\kappa(h)$
and is excluded afterwards.

(v) After (e1), $\Used(u)$ consists of one junction for each coloured item of a
non-ultra hub at $u$. The elements of $E'(u)$ correspond to further live items at
$u$: a PAR end at $u$ corresponds to the $\Jpar$ item or the leg of the cherry at
$u$, and an ultra hub end is itself an item. So
$|\Used(u)|+|E'(u)|\le M_l-1$. Every port carrying a live item is JV-good, so
$|\Cand_l(u)\setminus\Used(u)|\ge\Hcd_l-(M_l-1)\ge\Hcd_l/2\ge M_l>|E'(u)|$,
using $\Hcd_l\ge2^{10}M_l^{10}\ge2M_l$.

\emph{Remark on numerics (not used).} An exact evaluation of the union bound
$K^5\sum_sT_s$ with $m=M_l-1$ has its maximum, about $0.0502$, near $M_l=19$;
it tends to $9/256$ as $M_l\to\infty$, which is the limit of $K^5T_4$, where
$T_4=\binom{m}{4}/\binom{K}{3}^3\sim9m^4/K^9$. An earlier numerical claim of $0.040$ for this quantity was
false. None of these numerical values is used: only the analytic bound (iii) is
used, and it holds since $M_l\ge2^{40}$.
\end{proof}

%% ----------------------------------------------------------------------
\subsection{The quotient is simple, and every decomposition lifts}

For a HUB colour $\kappa$, a hub $h$ and $w\in\Pool_l$, let $m_\kappa(h,w)$ be
the number of edges $h[w]$ of $\BH_\kappa$, that is, the number of coloured hub
items $(h,u)$ of colour $\kappa$ with junction $w$.

\begin{lemma}[Lemma MULT]\label{s7:lemMULT}
Let $3\le l\le R$, $\kappa\in[4M_l]$, let $h$ be a hub and let
$w\in\Pool_{l,r}$. Then
\[
m_\kappa(h,w)\ \le\ \mult_r(w)\ \le\ \mu_r(w).
\]
Moreover, the number of sub-layers $B^{\mathrm{H},(\iota)}_\kappa$ ($\iota\ge1$)
in which $[w]$ is not isolated equals $\max_{h}m_\kappa(h,w)$, and the number of
those in which $h$ is not isolated equals $\max_{w}m_\kappa(h,w)$.
\deps{\ref{s7:consRound}, \ref{s6:lemJplus}, \ref{s7:defCand}, \ref{s7:defPool}, \ref{s2:propOV}}
\end{lemma}

\begin{proof}
Let $(h,u_1),\dots,(h,u_m)$ be the coloured hub items of colour $\kappa$ at $h$
with junction $w$, so $m=m_\kappa(h,w)$. By step (e) of
Construction~\ref{s7:consRound}, the junction of $(h,u_i)$ lies in
$\Cand_l(u_i)\subseteq\Pool_{l,r(u_i)}$ (Definition~\ref{s7:defCand}). Every
vertex has a single pool label (Definition~\ref{s7:defPool}), and
$w\in\Pool_{l,r}$, so $r(u_i)=r$ for all $i$. Moreover
$w\in\Cand_l(u_i)\subseteq N_{H_{Y(u_i)}}(u_i)\subseteq V(Y(u_i))$; this is where
it matters that candidates are taken in the class graph $H_{Y(u_i)}$ itself. So
each $Y(u_i)$ is an ancestor of round $r$ whose vertex set contains $w$. The
items $(h,u_i)$ are distinct $\Jhub$ edges at $h$ of round $l$. By \Jp{1} of
Lemma~\ref{s6:lemJplus}, for each class $Y$ there is at most one $\Jhub$ edge of
class $Y$ at $h$ in round $l$, aggregated over all MIX-parts. So the classes
$Y(u_1),\dots,Y(u_m)$ are pairwise distinct. Hence $m\le\mult_r(w)$, the number
of ancestors of round $r$ containing $w$, and $\mult_r(w)\le\mu_r(w)$ by
Proposition~\ref{s2:propOV}.

For the second statement: $[w]$ is not isolated in $B^{\mathrm{H},(\iota)}_\kappa$
if and only if some edge of rank $\iota$ is incident to $[w]$. By the definition
of ranks, this holds if and only if some pair $\{h,[w]\}$ carries at least
$\iota$ edges, that is, $\max_h m_\kappa(h,w)\ge\iota$. The statement for $h$ is
proved in the same way.
\end{proof}

\begin{lemma}[$Q_l$ is simple]\label{s7:lemSimple}
Every sub-layer of Construction~\ref{s7:consRound}(g) is a simple graph: it has
no parallel edges, PAR sub-layers have no loops, and HUB sub-layers are
bipartite between hubs and junction copies $[w]$, where the hub $h$ and the
junction $w$ of any edge $h[w]$ are distinct vertices of $G$. Consequently $Q_l$
is a simple graph without isolated vertices.
\deps{\ref{s7:consRound}}
\end{lemma}

\begin{proof}
By the definition of ranks, a sub-layer contains at most one edge of rank
$\iota$ between any two vertices, so it has no parallel edges. An edge of a PAR
layer joins $[w_1]$ and $[w_2]$ with $w_1\ne w_2$, because looped PAR objects
were paid in (e3); so PAR sub-layers have no loops. An edge $h[w]$ of a HUB
layer joins the side $D_l$ to the side $\{[w]\}$. Moreover $h$ is the hub of a
coloured, hence live, item, so $h$ is a live vertex and $h\notin\Pool_l$ by (a2),
while $w\in\Pool_l$. So $h\ne w$ as vertices of $G$, and HUB sub-layers have no
loops even when $[w]$ is identified with $w$. So every sub-layer is simple.
$Q_l$ is a vertex-disjoint union of the sub-layers with their isolated vertices
deleted, so it is simple and has no isolated vertices.
\end{proof}

\begin{lemma}[Lemma JV-L: the lift]\label{s7:lemLift}
Let $3\le l\le R$, fix any past, any $\xi_l$, and let $\Live$ be the set of live
vertices of round $l$.
\begin{enumerate}
\item[(i)] $\Live\cap\Pool_l=\emptyset$. Every port carries at most one PAR
object of each PAR colour and at most one coloured hub item of each HUB colour.
Hence, although ports are not vertices of $Q_l$, at most one edge of each
sub-layer is an object or item with an end at a given port. Every fresh centre is the middle of at
most one cherry of each PAR colour. Hubs, fresh centres and ports are pairwise
disjoint.
\item[(ii)] For \emph{every} decomposition $\Dec$ of $E(Q_l)$ into cycles and
single edges, the lift of each cycle of $\Dec$ (Construction~\ref{s7:consRound}(h))
is a cycle of $G$, and the lift of each single edge of $\Dec$ consists of one or
two single edges of $G$. The lifts of distinct members of $\Dec$ are pairwise
edge-disjoint. Hence the lift of $\Dec$ consists of at most $2|\Dec|$ pairwise
edge-disjoint objects of $G$.
\item[(iii)] Every item is either paid or lies in exactly one edge of $Q_l$.
Every junction edge $uw$ is the junction edge of at most one item-end, and that
end is at $u$, never at $w$. The edge $uw$ lies in exactly one lent class, namely
$\LJV_{Y(u),l}$, where $Y(u)$ is a lend-good ancestor of round at most $l-2$; the
only consumer of this class is the round-$l$ step. Junction edges lie in
$E_{r}(Y)$ for ancestors $Y$ of rounds $r\le l-2$, items lie in $E_l(Z)$ for
$Z\in\Std_l$, and these edge sets are disjoint.
\item[(iv)] Consequently the round step of Construction~\ref{s7:consRound},
applied at every round $3\le l\le R$, is a J-consumer:
$\Obj_l$ and $\LentJV_l$ satisfy \JC{1}--\JC{3} of
Definition~\ref{s6:defJconsumer}. The objects of $\Obj_l$ other than paid single
edges number at most $2|\Dec_l|=2\f(Q_l)$.
\end{enumerate}
\deps{\ref{s7:consRound}, \ref{s7:lemSimple}, \ref{s7:lemWellDef}, \ref{s6:lemJplus}, \ref{s6:defJconsumer}, \ref{s3:defCOL}, \ref{s2:propStructure}, \ref{s6:defLending}}
\fixes{\RTa{M1}, \RTb{J2}}
\end{lemma}

\begin{proof}
(i) The first claim is (a2) of Construction~\ref{s7:consRound}. The claims about
ports and fresh centres are Lemma~\ref{s7:lemWellDef}(i) together with the SDR
of step (d). Since every edge of a sub-layer of colour $\kappa$ is a PAR object or
a coloured hub item of colour $\kappa$, at most one edge of each sub-layer is an
object or item with an end at a given port. Disjointness of the three roles is \Jp{3} of
Lemma~\ref{s6:lemJplus}.

(ii) \emph{PAR cycles.} Let $[w_0]\cdots[w_{m-1}]$ be a cycle of a PAR sub-layer
of colour $\kappa$. By Lemma~\ref{s7:lemSimple} the sub-layer is simple, so
$m\ge3$. Distinct vertices of the sub-layer are copies of distinct pooled
vertices, so $w_0,\dots,w_{m-1}$ are distinct. The objects $o_0,\dots,o_{m-1}$
are distinct, because they are distinct edges of the cycle, and all have colour
$\kappa$. By (i) a port carries at most one object of colour $\kappa$, so ports
of distinct objects are distinct. The two ends of one PAR object lie at distinct
ports ($u\ne v$ for a $\Jpar$ item, $u\ne u'$ for a cherry), so $p_i\ne q_i$.
Hence the $2m$ ports $p_i,q_i$ are pairwise distinct. The middles $x_i$ of the
cherries among the $o_i$ are distinct, because a fresh centre is the middle of at
most one cherry of colour $\kappa$. Ports lie in $\Qs_Z$ and middles in $F_Z$, so
by (i) no port is a middle. Ports and middles are live, junctions are pooled,
and $\Live\cap\Pool_l=\emptyset$. So all vertices of the closed walk
$w_0p_0(x_0)q_0w_1\cdots q_{m-1}w_0$ are pairwise distinct. Consecutive vertices
are adjacent in $G$: $w_ip_i$ and $q_iw_{i+1}$ are junction edges (edges of
$H_{Y(p_i)}$ and $H_{Y(q_i)}$), and $p_iq_i$ or $p_ix_i,x_iq_i$ are the J-edges
of $o_i$. The walk has length at least $3m\ge9$. A closed walk of length at
least $3$ with pairwise distinct vertices is a cycle of $G$.

\emph{HUB cycles.} Let $h_0[w_0]h_1[w_1]\cdots h_{m-1}[w_{m-1}]h_0$ be a cycle of
a HUB sub-layer of colour $\kappa$. The sub-layer is simple and bipartite
(Lemma~\ref{s7:lemSimple}), so the cycle has even length at least $4$, and
$m\ge2$. The hubs $h_i$ are pairwise distinct, and so are the junctions $w_i$.
The $2m$ items $(h_i,u_i)$ and $(h_{i+1},u'_i)$ are distinct, because they are
distinct edges of the cycle, and all have colour $\kappa$. A port carries at most
one coloured hub item of colour $\kappa$, so the $2m$ ports $u_i,u'_i$ are
pairwise distinct; in particular $u_i\ne u'_i$. Hubs and ports are disjoint by
(i), and hubs and ports are live while junctions are pooled. So all vertices of
the closed walk $h_0u_0w_0u'_0h_1\cdots u'_{m-1}h_0$ are distinct. Consecutive
vertices are adjacent: $h_iu_i$ and $u'_ih_{i+1}$ are items, and $u_iw_i$,
$w_iu'_i$ are junction edges. The length is $4m\ge8$, so the lift is a cycle of
$G$.

\emph{Single edges.} A single edge of a PAR sub-layer is lifted to the one or two
J-edges of its PAR object, and a single edge of a HUB sub-layer to one item. So
every member of $\Dec$ lifts to at most two objects.

\emph{Edge-disjointness.} Every live item that is not paid is a coloured hub
item or belongs to exactly one unpaid PAR object, so it lies in exactly one edge
of $Q_l$. Each edge of $Q_l$ lies in exactly one member of $\Dec$. So the item
edges used by the lifts of distinct members are distinct. A junction edge $uw$
on a lifted cycle is the junction edge of an end at $u$ with junction $w$. The
ends at $u$ have pairwise distinct junctions: in (e1) every chosen junction is
added to $\Used(u)$ and is excluded later, and in (e2) the map is injective into
$\Cand_l(u)\setminus\Used(u)$. So $uw$ determines the end at $u$, hence the
item, the edge of $Q_l$ and the member of $\Dec$. It cannot also be the junction
edge of an end at $w$: $w$ is pooled, so every item with vertex $w$ was paid in
(a2), and no end lies at $w$. Finally, junction edges are not items: a junction
edge lies in $E(H_Y)\subseteq E_r(Y)$ for an ancestor $Y$ of round $r\le l-2$,
an item lies in $E_l(Z)$ with $Z\in\Std_l$, and these edge sets are disjoint
(Proposition~\ref{s2:propStructure}(iii); Lemma~\ref{s6:lemJplus}(iii)).

(iii) The first two claims were shown in the proof of (ii); every item that is
not in an edge of $Q_l$ was paid in (a2), (a3), (b), (d) or (e3). Let $uw$ be a
junction edge of an end at $u$. Then $w\in\Cand_l(u)$, so
$uw\in\LJV_{Y(u),l}\subseteq E(H_{Y(u)})$. The graphs $H_Y$ of distinct ancestors
are edge-disjoint, since $E(H_Y)\subseteq E_r(Y)$ and these sets are pairwise
disjoint (Proposition~\ref{s2:propStructure}(iii)). Inside $\Lend_{Y(u)}$ the
edge has exactly one colour. So $uw$ lies in exactly one lent class, namely
$\LJV_{Y(u),l}$. By COL(g) of Definition~\ref{s3:defCOL}, the only consumer of
this class is the round-$l$ J-consumer. The port $u$ carries an end, so it is
live, and hence $u\in\Qs_Z=Q_Z\setminus\Lost_Z$. Since $\Lost_Z$ contains every
classed port whose class is lend-bad, $Y(u)$ is lend-good. Finally
$r(u)\le l-2$. The disjointness claim was shown in (ii).

(iv) \JC{1}: $\LentJV_l$ consists of junction edges on lifted cycles, and by
(iii) each lies in $\LJV_{Y,l}$ for a lend-good ancestor $Y$ of round at most
$l-2$. \JC{2}: the paid single edges are J-edges that lie in no edge of $Q_l$;
they are distinct from each other and from all edges used by the lifts, which
consist of items in edges of $Q_l$ and of junction edges. Together with (ii), the
objects of $\Obj_l$ are pairwise edge-disjoint cycles and single edges. Their
union contains every edge of $\Jlost_l$ (paid in (a1)) and every item: each item
is paid or lies in an edge of $Q_l$, whose member of $\Dec_l$ lifts to objects
containing all items of that edge. The only other edges used are the junction
edges on lifted cycles, which form $\LentJV_l$ by definition. So the union is
exactly $J_l\cup\LentJV_l$. \JC{3} holds because no other edge is used. The
bound $2|\Dec_l|=2\f(Q_l)$ is (ii) for $\Dec=\Dec_l$. Junction edges that were
assigned but do not lie on a lifted cycle are not in $\LentJV_l$; they stay in
$E_r(Y)\setminus\Lent(Y)$ as junk for their owner.
\end{proof}

%% ----------------------------------------------------------------------
\subsection{Quotient size and payments}

In this subsection $3\le l\le R$ and the past $\Past_l$ are fixed. We say that
\emph{the lists are fixed} when, in addition, the variables $\eta_h$ and
$\zeta_{h,u}$ of $\xi_l$ are fixed. Then steps (a)--(d) and (e1) of
Construction~\ref{s7:consRound}, the sets $E'(u)$ and their orders
$e_1,\dots,e_q$ are determined (by the restrictions on the fixed rules), and
only the orders $\prec_u$ are random. By step (e2), the junctions of the ends at distinct ports are then
independent, and the junction of an end in $E'(u)$ is uniform on
$\Cand_l(u)\setminus\Used(u)$, a set of size at least $\Hcd_l/2$
(Lemma~\ref{s7:lemWellDef}(v)). By \Jp{2} of Lemma~\ref{s6:lemJplus},
$|J_l|\le n(M_l-1)\le1.37n(M_l-1)$. In this subsection we deliberately use the
looser second form (it only enlarges constants). Hence the number of PAR
objects and the number of hub items are each at most $1.37nM_l$. For a PAR colour $\kappa$ and $w\ne w'$ in
$\Pool_l$, let $m_\kappa(w,w')$ be the number of edges $[w][w']$ of $\BP_\kappa$.
As in Lemma~\ref{s7:lemMULT}, $[w]$ is not isolated in exactly
$\max_{w'}m_\kappa(w,w')$ of the sub-layers $B^{\mathrm{P},(\iota)}_\kappa$.

\begin{lemma}[Lemma CC: PAR copies]\label{s7:lemCC}
Let the past and the lists be fixed. Fix a PAR colour $\kappa$ and $w\in\Pool_l$.
Condition further on the indicators
\[
I_x\defeq\ind{\text{the end at $x$ of the $\kappa$-object at $x$ receives the junction }w}
\]
for every port $x$ carrying a PAR object of colour $\kappa$. Let $S_w$ be the
set of PAR objects of colour $\kappa$ with exactly one end whose junction is $w$.
\begin{enumerate}
\item[(i)] The partner ports $v_o$ of the objects $o\in S_w$ (the ports of the
ends whose junction is not $w$) are pairwise distinct. Under the conditioning,
the junctions of the partner ends are independent, and the junction of the
partner end of $o$ is uniform on $(\Cand_l(v_o)\setminus\Used(v_o))\setminus\{w\}$.
In particular it equals a given $w'$ with probability at most $3/\Hcd_l$.
\item[(ii)] For every integer $t\ge1$,
\[
\Ex\Bigl[\max_{w'}m_\kappa(w,w')\ \Big|\ \Past_l,\text{lists},(I_x)_x\Bigr]
\ \le\ t\cdot\ind{S_w\ne\emptyset}+\frac{6e|S_w|}{\Hcd_l}+4e\,2^{-t}|S_w| .
\]
\item[(iii)] With $t\defeq\tCC_l\defeq\lceil2\log_2M_l\rceil$, the number of
vertices of $Q_l$ lying in PAR sub-layers satisfies
\[
\Ex\bigl[\#\{\text{PAR copies}\}\ \big|\ \Past_l,\text{lists}\bigr]\ \le\
3M_l(\tCC_l+1)|\Pool_l|+44.7\,\frac{nM_l}{\Hcd_l}+29.8\,\frac{n}{M_l}.
\]
\end{enumerate}
These bounds hold for every past, in particular for every realized $J_l$, and for
every outcome of the lists, the SDR and the greedy step (e1): only the orders
$\prec_u$ are random.
\deps{\ref{s7:consRound}, \ref{s7:lemWellDef}, \ref{s1:citChernoffGen}, \ref{s6:lemJplus}}
\fixes{\RTb{I4}, \RTa{M2}}
\end{lemma}

\begin{proof}
The edges at $[w]$ in $\BP_\kappa$ are exactly the objects in $S_w$. Indeed, an
object with both ends at junction $w$ is looped and was paid in (e3), and an
object with no end at junction $w$ is not incident to $[w]$. So
$m_\kappa(w,w')$ is the number of $o\in S_w$ whose partner end has junction
$w'$.

(i) A PAR object has its two ends at distinct ports, and every port carries at
most one PAR object of colour $\kappa$ (Lemma~\ref{s7:lemWellDef}(i)). So every
such port $x$ carries exactly one end of colour $\kappa$, the indicator $I_x$ is
well defined, and the partner ports of distinct objects of $S_w$ are distinct.
The set $S_w$ and the partner ports are determined by the indicators. With the
past and the lists fixed, $\Cand_l(x)\setminus\Used(x)$, $E'(x)$ and its order
are fixed for every port $x$ (restrictions on the fixed rules in
Construction~\ref{s7:consRound}), so the junction of the $\kappa$-end at a
port $x$ is a function of $\prec_x$ alone, and the orders of distinct ports are independent.
So conditioning on $(I_x)_x$ keeps the partner ends independent, and it
conditions the partner end at $v_o$ only on the event $I_{v_o}=0$. The junction
of that end is uniform on $C(v_o)\defeq\Cand_l(v_o)\setminus\Used(v_o)$, since
it is the image of a fixed element under a uniform injection. Given $I_{v_o}=0$
it is therefore uniform on $C(v_o)\setminus\{w\}$. By
Lemma~\ref{s7:lemWellDef}(v), $|C(v_o)\setminus\{w\}|\ge\Hcd_l/2-1\ge\Hcd_l/3$.

(ii) If $S_w=\emptyset$ the left side is $0$. Otherwise, by (i), each
$m_\kappa(w,w')$ is, under the conditioning, a sum of independent indicators
with mean $\mu_{w'}\le3|S_w|/\Hcd_l$, and $\sum_{w'}\mu_{w'}=|S_w|$. Put
$T\defeq\max(t,\lceil6e|S_w|/\Hcd_l\rceil)$. For every integer $j\ge T$ we have
$j\ge2e\mu_{w'}$. By Cited result~\ref{s1:citChernoffGen},
$\Prob(m_\kappa(w,w')\ge j)\le(e\mu_{w'}/j)^j\le(e\mu_{w'}/j)\,2^{1-j}$. Hence,
writing $m$ for $m_\kappa(w,w')$,
\begin{align*}
\Ex\bigl[m\,\ind{m\ge T}\bigr]
&=T\,\Prob(m\ge T)+\sum_{j>T}\Prob(m\ge j)
\le e\mu_{w'}2^{1-T}+\sum_{j>T}\frac{e\mu_{w'}}{j}2^{1-j}\\
&\le e\mu_{w'}\,2^{2-T}\le e\mu_{w'}\,2^{2-t}.
\end{align*}
Moreover
\[
\max_{w'}m_\kappa(w,w')\le(T-1)+\sum_{w'}m_\kappa(w,w')\,\ind{m_\kappa(w,w')\ge T}:
\]
if the maximum is below $T$ it is at most $T-1$, and otherwise it is one of the
summands. Taking conditional expectations,
\[
\Ex\Bigl[\max_{w'}m_\kappa(w,w')\ \Big|\ \cdot\Bigr]\le T-1+4e\,2^{-t}|S_w|
\le t+\frac{6e|S_w|}{\Hcd_l}+4e\,2^{-t}|S_w|,
\]
because $T-1\le\max(t-1,6e|S_w|/\Hcd_l)$.

(iii) The number of PAR copies is
$\sum_{\kappa\in[3M_l]}\sum_{w\in\Pool_l}\max_{w'}m_\kappa(w,w')$. Averaging (ii)
over the indicators and summing over the at most $3M_l|\Pool_l|$ pairs
$(\kappa,w)$ gives at most
$3M_l\,t\,|\Pool_l|+\bigl(6e/\Hcd_l+4e\,2^{-t}\bigr)\,\Ex\sum_{\kappa,w}|S_w|$.
Every PAR object lies in at most two of the sets $S_w$ (one for each end), so
$\sum_{\kappa,w}|S_w|\le2\cdot1.37nM_l=2.74\,nM_l$. With $t=\tCC_l$ we have
$2^{-t}\le M_l^{-2}$, and $6e\cdot2.74<44.7$ and $4e\cdot2.74<29.8$. Finally
$t\le t+1$.
\end{proof}

\begin{lemma}[Ultra-hub copies]\label{s7:lemUltra}
Let the past and the lists be fixed.
\begin{enumerate}
\item[(i)] Let $h$ be an ultra hub, $\kappa\in[4M_l]$, and let
$\omega_1,\dots,\omega_N$ be the junctions of the coloured hub items of $h$ of
colour $\kappa$. These items lie at $N$ distinct ports, because distinct items
at $h$ are distinct edges $hu$ and $G$ is simple. The $\omega_i$ are independent,
and each is uniform on a set of size at least $\Hcd_l/2$.
\item[(ii)] If $N\ge1$ then, with $\mu^*\defeq2N/\Hcd_l$,
$\Ex[\max_w m_\kappa(h,w)\mid\Past_l,\text{lists}]\le\log_2N+2e\mu^*+8$.
\item[(iii)] The number of vertices of HUB sub-layers that are copies of $h$
satisfies
$\Ex[\,\cdot\mid\Past_l,\text{lists}]\le4M_l(\log_2\clive_h+8)+4e\,\clive_h/\Hcd_l$.
\item[(iv)] Summed over all ultra hubs, the expected number of hub copies of
ultra hubs, given $\Past_l$ and the lists, is at most
\[
320\,\frac{nM_l^{3}\,(\log_2\thult_l+8)}{\lambda_{l-2}^{95}} .
\]
\end{enumerate}
The bound depends on the past only through deterministic counts, so it holds for
every past and every realized $J_l$.
\deps{\ref{s7:consRound}, \ref{s7:lemWellDef}, \ref{s6:lemJplus}, \ref{s1:citChernoffGen}, \ref{s7:lemMULT}, \ref{s7:lemCand}}
\fixes{\RTa{M8}}
\end{lemma}

\begin{proof}
(i) Items of ultra hubs receive their junctions in (e2). The junction of the
item $(h,u_i)$ is the image of its end under the uniformly random injection at
$u_i$, so it is uniform on $\Cand_l(u_i)\setminus\Used(u_i)$, a set of size at
least $\Hcd_l/2$ by Lemma~\ref{s7:lemWellDef}(v). The ports $u_i$ are distinct,
and injections at distinct ports are independent, so the $\omega_i$ are
independent. (Distinctness of the ports comes from the simplicity of $G$; it
does not use \Jp{1}.)

(ii) Let $m_w\defeq\#\{i:\omega_i=w\}$. It is a sum of independent indicators
with mean $\mu_w\le2N/\Hcd_l=\mu^*$, and $\sum_w\mu_w=N$. Put
$T\defeq\max(\lceil2e\mu^*\rceil,\lceil\log_2N\rceil+1)$. For $j\ge T$ we have
$e\mu_w/j\le1/2$, so by Cited result~\ref{s1:citChernoffGen}
$\Prob(m_w\ge j)\le(e\mu_w/j)^j\le(e\mu_w/j)2^{1-j}$, and summing over $w$ gives
$\sum_w\Prob(m_w\ge j)\le(eN/j)2^{1-j}$. Hence
\[
\Ex\bigl[\max_wm_w\bigr]=\sum_{j\ge1}\Prob\bigl(\max_wm_w\ge j\bigr)
\le(T-1)+\sum_{j\ge T}\frac{eN}{j}2^{1-j}\le(T-1)+\frac{eN}{T}2^{2-T}
\le T-1+\frac{2e}{T},
\]
using $2^{-T}\le2^{-\log_2N-1}=1/(2N)$. Finally
$T-1\le\max(2e\mu^*,\log_2N+1)\le\log_2N+2e\mu^*+1$ and $2e/T\le2e<7$.

(iii) By Lemma~\ref{s7:lemMULT}, the number of copies of $h$ in the sub-layers of
colour $\kappa$ is $\max_wm_\kappa(h,w)$. Let $N_\kappa$ be the number of
coloured items of $h$ of colour $\kappa$. Then $N_\kappa\le\clive_h$ and
$\sum_\kappa N_\kappa\le\clive_h$, and at most $4M_l$ colours have
$N_\kappa\ge1$. By (ii) the expected number of copies of $h$ is at most
$\sum_{\kappa:N_\kappa\ge1}\bigl(\log_2\clive_h+8\bigr)+2e\sum_\kappa2N_\kappa/\Hcd_l
\le4M_l(\log_2\clive_h+8)+4e\,\clive_h/\Hcd_l$.

(iv) The function $g(x)\defeq(\log_2x+8)/x$ is decreasing on $[1,\infty)$, since
$x^2g'(x)=1/\ln2-8-\log_2x<0$ there. For an ultra hub, $\clive_h>\thult_l\ge1$,
so $\log_2\clive_h+8=\clive_h\,g(\clive_h)\le\clive_h(\log_2\thult_l+8)/\thult_l$.
By \Jp{2} of Lemma~\ref{s6:lemJplus},
$\sum_h\clive_h\le|J_l|\le1.37nM_l$. We have
$M_l\Hcd_l/7=\lambda_{l-2}^{95}/(56M_l)$, and $M_l\Hcd_l\ge2^{10}M_l^{11}$
(Lemma~\ref{s7:lemCand}(iv)), so $\lambda_{l-2}^{95}/M_l\ge2^{453}$. Hence
$\thult_l\ge\lambda_{l-2}^{95}/(56M_l)-1\ge\lambda_{l-2}^{95}/(57M_l)$, and also
$\thult_l\ge2^{8}$. Summing (iii) over ultra hubs,
\[
\sum_{h\text{ ultra}}\Bigl(4M_l(\log_2\clive_h+8)+\frac{4e\,\clive_h}{\Hcd_l}\Bigr)
\le4M_l\cdot1.37nM_l\cdot\frac{57M_l}{\lambda_{l-2}^{95}}(\log_2\thult_l+8)
+4e\cdot1.37nM_l\cdot\frac{8M_l^2}{\lambda_{l-2}^{95}} .
\]
The first term is $312.36\,nM_l^3(\log_2\thult_l+8)/\lambda_{l-2}^{95}$ and the
second is at most $119.2\,nM_l^3/\lambda_{l-2}^{95}$. Since
$\log_2\thult_l+8\ge16$, the second term is at most
$7.5\,nM_l^3(\log_2\thult_l+8)/\lambda_{l-2}^{95}$, and $312.36+7.5<320$.
\end{proof}

\begin{lemma}[The stage-1 weight $\XV$]\label{s7:lemVstar}
For $3\le l\le R$ put
\[
X_{\mathrm{V},l}\defeq3M_l(\tCC_l+1)|\Pool_l|+4M_l\sum_{w\in\Pool_l}\mult_{r(w)}(w),
\qquad \XV\defeq\sum_{l=3}^{R}X_{\mathrm{V},l}.
\]
Each $X_{\mathrm{V},l}$ is a deterministic function of the run and the pool
labels, and
\[
\Ex X_{\mathrm{V},l}\le\frac{(6\log_2M_l+12)\,n}{M_l},\qquad
\Ex\XV\le2.1\,\frac{(6\log_2\Dstar+12)\,n}{\Dstar}.
\]
\deps{\ref{s7:defPool}, \ref{s2:propOV}, \ref{s2:lemTower}}
\fixes{\RTa{M3}, \RTa{M2}, \RTb{I4}}
\end{lemma}

\begin{proof}
The quantities $M_l$, $\tCC_l$ and $\mult_r(w)$ are determined by the run, and
$\Pool_l$ and $r(w)$ by the pool labels. Since $\tCC_l+1\le2\log_2M_l+2$ and
$\Ex|\Pool_l|\le q_ln=n/M_l^2$, the first summand has expectation at most
$(6\log_2M_l+6)n/M_l$. For the second,
\[
\Ex\sum_{w\in\Pool_l}\mult_{r(w)}(w)=\sum_{w}\sum_{r=1}^{l-2}q_l\pi_{l,r}\,\mult_r(w)
=q_l\sum_{r=1}^{l-2}\pi_{l,r}\sum_w\mult_r(w)\le q_l\cdot1.37n ,
\]
since $\sum_w\mult_r(w)\le1.37n$ for every $r$ (Proposition~\ref{s2:propOV}) and
$\sum_r\pi_{l,r}<1$. So the second summand has expectation at most
$4M_l\cdot1.37n/M_l^2=5.48\,n/M_l$. Adding, $\Ex X_{\mathrm{V},l}\le(6\log_2M_l+11.48)n/M_l$.
Summing over $l$ and using $\sum_l\psi(M_l)\le2.1\,\psi(\Dstar)$ for
$\psi(x)=(6\log_2x+12)/x$ (Lemma~\ref{s2:lemTower}(b)) gives the bound for
$\Ex\XV$.
\end{proof}

\begin{lemma}[Payments]\label{s7:lemPay}
For every past and every $3\le l\le R$ the following hold.
\begin{enumerate}
\item[(a1)] $|\Jlost_l|\le(M_l-1)|\Lost_l|$.
\item[(a2)] The number of items with a vertex in $\Pool_l$ is at most
$\sum_{v\in\Pool_l}\omega_l(v)$, where $\omega_l(v)\defeq\cagg_{v,l}$ if
$v\in D_l$ and $\omega_l(v)\defeq M_l$ otherwise.
\item[(a3)] The number of items with a JV-bad port end is at most $M_l-1$ times
the number of JV-bad classed ports of round $l$.
\item[(b)] Every fresh centre has at most one unpaired fresh leg per round, so
the number of unpaired fresh legs over all rounds is at most
$\sum_{l,Z}|F_Z|\le2n$.
\item[(c)] $\Ex[\text{SDR-failure payments}\mid\Past_l]\le n(M_l-1)(4M_l)^{-4}\le n/(256M_l^3)$,
and $\Ex[\text{loop-paid edges}\mid\Past_l]\le2\cdot(2/\Hcd_l)\cdot\#\{\text{PAR objects}\}\le5.5\,nM_l/\Hcd_l$.
\item[(d)] If $\xi_l$ lies in the event of Construction~\ref{s7:consRound}(f),
then $\payrd_l\le4\bigl(n/(256M_l^3)+5.5\,nM_l/\Hcd_l\bigr)\le2n/M_l$; if this
holds at every round, then $\sum_l\payrd_l\le9n/\Dstar$.
\item[(e)] No other item is paid. In particular, no item is paid because of a
coincidence of junctions at an ultra hub, and no item is paid for a concentrated
pair.
\end{enumerate}
\deps{\ref{s7:consRound}, \ref{s7:lemWellDef}, \ref{s6:lemJplus}, \ref{s2:propParentless}, \ref{s2:lemTower}}
\fixes{\RTb{J4}, \RTa{M8}}
\end{lemma}

\begin{proof}
We use the properties of Lemma~\ref{s6:lemJplus}.

(a1) Every $\Jlost$ edge joins a vertex of $\Lost_l$ to a port, and every vertex
carries at most $M_l-1$ J-edges by \Jp{2}.

(a2) Let $v\in\Pool_l$. If $v\in D_l$, then $v$ is a hub and not a port or a
fresh centre, by \Jp{3}. The items with vertex $v$ are therefore $\Jhub$ items
$(v,u)$: $\Jpar$ items have both ends in $\Qs_Z$, and $\Jfr$ items have their
centre in $F_Z$. By \Jp{1} there is at most one such item per class $Y$. Each
such class has $d_{Y,l}(v)\ge1$, because the item is an edge $vu\in E_l(Z)$ with
$u\in Q_Z$ and $Y(u)=Y$. So there are at most $\cagg_{v,l}$ items with vertex
$v$. If $v\notin D_l$, then $v$ carries at most $M_l-1\le M_l$ J-edges by \Jp{2}.

(a3) A port carries at most $M_l-1$ items by \Jp{2}.

(b) The pairing at a fresh centre $x$ leaves at most one live $\Jfr$ item
unpaired. Every fresh centre of round $l$ lies in $F_Z$ for some $Z\in\Std_l$,
and $\sum_{l,Z}|F_Z|\le2n$ by Proposition~\ref{s2:propParentless}(i).

(c) By \Jp{3}, every non-hub vertex has one role in one part per round, so there
are at most $n$ ports. A port $u$ pays at most $M_l-1$ items for an SDR failure,
and it fails with probability at most $(4M_l)^{-4}$ given $\Past_l$
(Lemma~\ref{s7:lemWellDef}(iii)). This gives the first bound. For loops, fix the
lists. A PAR object has its two ends at distinct ports $p\ne q$. Their junctions
are independent, and the junction at $q$ is uniform on a set of size at least
$\Hcd_l/2$ (Lemma~\ref{s7:lemWellDef}(v)). So the probability that the two
junctions coincide is
$\sum_w\Prob(p\to w)\Prob(q\to w)\le\max_w\Prob(q\to w)\le2/\Hcd_l$. A looped
object costs at most two edges (a cherry costs two, a $\Jpar$ item one), and
there are at most $1.37nM_l$ PAR objects (\Jp{2}). So the expected number of
loop-paid edges is at most $4\cdot1.37\,nM_l/\Hcd_l\le5.5\,nM_l/\Hcd_l$. This
bound does not depend on the lists, so it also holds given $\Past_l$ alone.

(d) The first inequality is the definition of the event in (f), together with
(c). Next, $4\cdot5.5\,nM_l/\Hcd_l=176\,nM_l^3/\lambda_{l-2}^{95}$. By
Lemma~\ref{s2:lemTower}(b), $M_l\le\lambda_{l-2}^{1.6}$, so
$\lambda_{l-2}\ge M_l^{1/1.6}\ge2^{25}$ and
$176\,M_l^3/\lambda_{l-2}^{95}\le176\,\lambda_{l-2}^{-90.2}\le\lambda_{l-2}^{-1.6}\le1/M_l$.
Also $4n/(256M_l^3)\le n/M_l$. Hence $\payrd_l\le2n/M_l$, and
$\sum_l\payrd_l\le2n\sum_l1/M_l\le4n/\Dstar\le9n/\Dstar$ by
Lemma~\ref{s2:lemTower}(b).

(e) In Construction~\ref{s7:consRound}, single edges are paid only in (a1),
(a2), (a3), (b), (d) and (e3).
\end{proof}

\begin{definition}[The stage-1 functional $\Xp$]\label{s7:defXprime}
With $\omega_l$ as in Lemma~\ref{s7:lemPay}(a2), put
\[
\Xpool\defeq\sum_{l=3}^{R}\Bigl[\sum_{v\in\Pool_l}\omega_l(v)+(M_l-1)\cdot\#\{\text{JV-bad classed ports of round }l\}\Bigr],
\]
let $\XV$ be as in Lemma~\ref{s7:lemVstar}, and let $\XU$ be as in
Definition~\ref{s6:defLending}. Then
\[
\Xp\defeq\XU+\Xpool+\XV .
\]
Each of the three terms is non-negative and is a deterministic function of the
run, $\delta$ and the stage-1 outcome. Indeed, $\omega_l(v)$ is determined by the
run and $\delta$ (through $\cagg_{v,l}$), $\Pool_l$ is determined by the pool
labels, and JV-badness is determined by stage~1. $\Xp$ is the functional used in
stage (1e) of the schedule.
\deps{\ref{s6:defLending}, \ref{s7:lemVstar}, \ref{s7:lemPay}}
\end{definition}

\begin{lemma}[Expectation of $\Xp$]\label{s7:lemEXprime}
$\Ex\Xp\le\epsX(\Dstar)\,n$, where
\[
\epsX(\Dstar)\defeq\epsU(\Dstar)+\frac{10.3}{\Dstar}+2.1\,\frac{6\log_2\Dstar+12}{\Dstar}
\]
and $\epsU$ is the function of Lemma~\ref{s5:lemExpect} (through
Lemma~\ref{s6:lemLost}). More precisely, $\Ex\XU\le\epsU(\Dstar)n+5.5n/\Dstar$,
$\Ex\Xpool\le4.8n/\Dstar$ and $\Ex\XV\le2.1(6\log_2\Dstar+12)n/\Dstar$.
\deps{\ref{s7:defXprime}, \ref{s6:lemLost}, \ref{s7:lemCand}, \ref{s7:lemVstar}, \ref{s2:propOV}, \ref{s2:propStructure}, \ref{s2:lemTower}, \ref{s5:lemExpect}, \ref{s2:lemCap}}
\fixes{\RTa{M5}}
\end{lemma}

\begin{proof}
The bound on $\Ex\XU$ is Lemma~\ref{s6:lemLost}, and the bound on $\Ex\XV$ is
Lemma~\ref{s7:lemVstar}. For $\Xpool$, fix a round $l$.

\emph{Pooled weights.} The weights $\omega_l(v)$ are determined by the run and
$\delta$, so they are independent of the pool labels. Hence
$\Ex\sum_{v\in\Pool_l}\omega_l(v)\le q_l\sum_v\omega_l(v)\le q_l\bigl(\sum_{h\in D_l}\cagg_{h,l}+nM_l\bigr)$.
Now $\cagg_{h,l}\le\sum_Yd_{Y,l}(h)$, and $\sum_Yd_{Y,l}(h)$ counts the edges
$hu\in E_l(Z)$ with $Z\in\Std_l$ and $u\in Q_Z$. Summing over $h$ counts each such
edge at most once for each of its classed-port ends. So
$\sum_h\cagg_{h,l}\le\sum_{Z\in\Std_l}\sum_{u\in Q_Z}\deg_{E_l(Z)}(u)$. Every
vertex has $E_l(Z)$-degree at most $M_l-1$ (Lemma~\ref{s2:lemCap}(ii)), and
$\sum_{Z\in\Std_l}|Z^0|\le1.37n$ by Proposition~\ref{s2:propOV}. Hence
$\sum_{h\in D_l}\cagg_{h,l}\le1.37nM_l$, and the expected pooled weight is at
most $2.37\,n/M_l$.

\emph{JV-bad ports.} There are at most $n$ classed ports of round $l$, because
$Q_Z\subseteq U_Z$ and the port sets $U_Z$ of distinct $Z\in\Std_l$ are pairwise
disjoint (Proposition~\ref{s2:propStructure}(iv)); this is the count used in
the proofs of \ref{s6:K6} of Lemma~\ref{s6:lemLost} and of \Jp{2} of
Lemma~\ref{s6:lemJplus}. By Lemma~\ref{s7:lemCand}(iii) each is JV-bad with
probability at most $\exp(-\Hcd_l/4)$, and $\exp(-\Hcd_l/4)\le M_l^{-3}$ because
$\Hcd_l\ge2^{10}M_l^{10}$ (Lemma~\ref{s7:lemCand}(iv)). So the expected JV-bad
weight is at most
\[
(M_l-1)\cdot n\cdot M_l^{-3}\ \le\ n/M_l^2\ \le\ 0.03\,n/M_l .
\]
The pooled-weight bound above counts the same classed ports by the
deliberately looser $\sum_{Z\in\Std_l}|Z^0|\le1.37n$; this only enlarges its
constant.

Hence $\Ex\Xpool\le\sum_l2.4\,n/M_l\le4.8\,n/\Dstar$, using
$\sum_l1/M_l\le2/\Dstar$ (Lemma~\ref{s2:lemTower}(b)). Adding the three bounds
gives $\Ex\Xp\le\epsU(\Dstar)n+(5.5+4.8)n/\Dstar+2.1(6\log_2\Dstar+12)n/\Dstar=\epsX(\Dstar)n$.
\end{proof}

\begin{lemma}[Lemma UH$^*$-split: quotient size]\label{s7:lemUHsplit}
For $3\le l\le R$ put
\[
\dett_l\defeq3|D_l|+78\,\frac{nM_l}{\Hcd_l}+30\,\frac{n}{M_l}
+320\,\frac{nM_l^{3}(\log_2\thult_l+8)}{\lambda_{l-2}^{95}},
\]
a deterministic function of the run.
\begin{enumerate}
\item[(i)] $X_{\mathrm{V},l}$ is the stage-1 weight of Lemma~\ref{s7:lemVstar}.
\item[(ii)] For every past, $\Ex[\copies_l\mid\Past_l]\le X_{\mathrm{V},l}+\dett_l$.
\item[(iii)] If the stage-1 outcome satisfies $\Xp\le3\,\Ex\Xp$ and at every round
$\xi_l$ lies in the event of Construction~\ref{s7:consRound}(f), then
\[
\sum_{l=3}^{R}|V(Q_l)|\ \le\ 12\,\Ex\Xp+4\sum_{l=3}^{R}\dett_l .
\]
\item[(iv)] Put $F(x)\defeq(3184+30400\log_2x)\,x^{-90.2}$ for $x\ge1$.
Then $\sum_l\dett_l\le3\epsA n+60n/\Dstar+2F(\log_2\Dstar)\,n$, and
$2F(\log_2\Dstar)<2^{-1000}$. Hence, in the situation of (iii),
\[
\sum_{l=3}^{R}|V(Q_l)|\le\thetaQ\,n,\qquad
\thetaQ\defeq\epsTwo(\Dstar)\defeq12\,\epsX(\Dstar)+4\Bigl(3\epsA+\frac{60}{\Dstar}+2F(\log_2\Dstar)\Bigr),
\]
with $\epsA=31\eps/(\Cp\log\log\Dstar)$ as in Lemma~\ref{s2:lemTower}(e). So
$\thetaQ$ depends on $\Dstar$ only, and $\thetaQ\le1/4$ by
condition~\ref{s1:condG4}.
\end{enumerate}
\deps{\ref{s7:lemCC}, \ref{s7:lemMULT}, \ref{s7:lemUltra}, \ref{s7:lemVstar}, \ref{s7:lemEXprime}, \ref{s7:defSchedule}, \ref{s2:propOV}, \ref{s2:lemTower}, \ref{s2:lemLacunary}, \ref{s1:condG4}, \ref{s1:condG1}, \ref{s7:consRound}, \ref{s7:lemWellDef}, \ref{s2:defHBtp}, \ref{s6:lemJplus}}
\fixes{\RTa{M3}, \RTa{M5}}
\end{lemma}

\begin{proof}
(ii) Fix the past and the lists. The vertices of $Q_l$ are of four kinds.

\emph{PAR copies.} Their conditional expectation is at most
$3M_l(\tCC_l+1)|\Pool_l|+44.7\,nM_l/\Hcd_l+29.8\,n/M_l$ by
Lemma~\ref{s7:lemCC}(iii).

\emph{Copies of non-ultra hubs} (deterministic given the lists). Let $h$ be
non-ultra. By Lemma~\ref{s7:lemWellDef}(iv) (see step (e1) of
Construction~\ref{s7:consRound}), $m_\kappa(h,w)\le1$ for all $\kappa$ and $w$.
So by Lemma~\ref{s7:lemMULT}, $h$ has exactly one copy in each colour that one of
its items uses, and none in the other colours. The colours of its items lie in
the union of its $k_h$ lists, which has $3k_h$ elements. So the number of copies
of non-ultra hubs is at most
$\sum_{h:\clive_h\ge1}3k_h\le\sum_{h\in D_l}3(1+8\clive_h/\Hcd_l)\le3|D_l|+24\cdot1.37\,nM_l/\Hcd_l$,
since $\sum_h\clive_h\le1.37nM_l$. Here $24\cdot1.37=32.88<32.9$.

\emph{Junction copies in HUB sub-layers} (the count depends on the junctions,
but the bound below is deterministic). By
Lemma~\ref{s7:lemMULT}, in colour $\kappa$ the vertex $[w]$ has
$\max_hm_\kappa(h,w)\le\mult_{r(w)}(w)$ copies. Summing over the $4M_l$ HUB
colours and $w\in\Pool_l$ gives at most $4M_l\sum_{w\in\Pool_l}\mult_{r(w)}(w)$.

\emph{Copies of ultra hubs.} Their conditional expectation is at most
$320\,nM_l^3(\log_2\thult_l+8)/\lambda_{l-2}^{95}$ by Lemma~\ref{s7:lemUltra}(iv).

The first summands of the PAR bound and of the junction bound add up to
$X_{\mathrm{V},l}$. The remaining terms add up to at most $\dett_l$, because
$44.7+32.9\le78$ and $29.8\le30$. Averaging over the lists gives (ii).

(iii) By the choice in (f) and by (ii),
$|V(Q_l)|=\copies_l\le4\,\Ex[\copies_l\mid\Past_l]\le4(X_{\mathrm{V},l}+\dett_l)$
at every round. Summing over $l$,
$\sum_l|V(Q_l)|\le4\XV+4\sum_l\dett_l\le4\Xp+4\sum_l\dett_l\le12\,\Ex\Xp+4\sum_l\dett_l$,
because $\XU,\Xpool\ge0$ and $\Xp\le3\Ex\Xp$ (stage (1e) of
Definition~\ref{s7:defSchedule}).

(iv) By Proposition~\ref{s2:propOV}(K2), $|D_l|\le7.6\eps n/\log P_l$. So
$\sum_l3|D_l|\le3n\sum_l7.6\eps/\log P_l\le1.5\epsA n\le3\epsA n$ by
Lemma~\ref{s2:lemTower}(e) (the factor~$2$ of slack is deliberate; it only
enlarges $\epsTwo$). Next,
$\sum_l30n/M_l\le60n/\Dstar$ by Lemma~\ref{s2:lemTower}(b). For the remaining
terms, $78\,nM_l/\Hcd_l=624\,nM_l^3/\lambda_{l-2}^{95}$, and
$\log_2\thult_l\le\log_2(M_l\Hcd_l/7)\le95\log_2\lambda_{l-2}$. With
$M_l\le\lambda_{l-2}^{1.6}$ (Lemma~\ref{s2:lemTower}(b)), these terms are at most
$nF(\lambda_{l-2})$, with $F$ as in (iv).
Here $\lambda_{l-2}\ge M_l^{1/1.6}\ge2^{25}$. On $[2^{25},\infty)$, $F$ is
decreasing, and $F(2x)\le F(x)/2$, since the first factor at most doubles while
$x^{-90.2}$ is multiplied by $2^{-90.2}$. By Lemma~\ref{s2:lemLacunary}(ii),
$\lambda_{r}\ge2^{\lambda_{r+1}/A}\ge2\lambda_{r+1}$ for consecutive rounds,
because $\lambda_{r+1}\ge2^{25}$. So $X_l\defeq F(\lambda_{l-2})$
($3\le l\le R$) satisfies $X_l\le X_{l+1}/2$. By
Lemma~\ref{s2:lemLacunary}(i), $\sum_lX_l\le2X_R=2F(\lambda_{R-2})$ (if $R\le2$
the sum is empty). Moreover $\lambda_{R-2}\ge\lambda_R=\log_2d_R\ge\log_2\Dstar$,
by Lemma~\ref{s2:lemTower}(a) and because round $R$ has $d_R\ge\Dstar$ by
rule~\ref{s2:R0} of Definition~\ref{s2:defHBtp}; and $\log_2\Dstar\ge2^{256}\ge2^{25}$ by \Gc{1}(a) at $\mu=\log_2\log_2\Dstar$. Since $F$
is decreasing on $[2^{25},\infty)$, we get
\[
\sum_lX_l\ \le\ 2F(\log_2\Dstar)\ \le\ 2F(2^{25})\ \le\ 2\cdot2^{20}\cdot2^{-2255}\ <\ 2^{-1000}.
\]
Together, $\sum_l\dett_l\le3\epsA n+60n/\Dstar+2F(\log_2\Dstar)\,n$. The bound
on $\sum_l|V(Q_l)|$ follows from (iii) and $\Ex\Xp\le\epsX(\Dstar)n$
(Lemma~\ref{s7:lemEXprime}). The functions $\epsX$, $\epsA$ and
$F(\log_2\Dstar)$ depend on $\Dstar$ only, and $\thetaQ\le1/4$ is part of
condition~\ref{s1:condG4}.
\end{proof}

%% ----------------------------------------------------------------------
\subsection{One joint outcome, the cost, and Theorem \texorpdfstring{JV$^{+*}$}{JV+*}}

\begin{lemma}[One joint outcome exists]\label{s7:lemOneOutcome}
There is a joint outcome, consisting of a stage-1 outcome, stage-3 outcomes for
all parts, and draws $\xi_R,\dots,\xi_3$, such that:
\begin{enumerate}
\item[(1)] $\Xp\le3\,\Ex\Xp$;
\item[(2)] the stage-3 labels of every part lie in that part's good event: the
event of Lemma~\ref{s4:lemTPV} for every standalone pre-part (with the data of
step (1) of Construction~\ref{s6:consOrder}); the event of Lemma~\ref{s4:lemPV}
as used in Lemma~\ref{s5:lemChild} for every non-demoted light part; and the event
of Theorem~\ref{s4:thmVXp} as used in Lemma~\ref{s5:lemDemoted} for every
demoted light part;
\item[(3)] at every round $3\le l\le R$, $\xi_l$ lies in the event of
Construction~\ref{s7:consRound}(f). Consequently
$\copies_l\le4(X_{\mathrm{V},l}+\dett_l)$ and
$\payrd_l\le4\bigl(n/(256M_l^3)+5.5\,nM_l/\Hcd_l\bigr)$.
\end{enumerate}
Such an outcome is obtained as follows.
(i) Choose the stage-1 outcome in $\{\Xp\le3\Ex\Xp\}$; stage~2 is then
determined.
(ii) For every part, choose its stage-3 labels in its good event.
(iii) For $l=R,R-1,\dots,3$ in turn: run steps (1)--(3) of
Construction~\ref{s6:consOrder} at round $l$, which are deterministic given the
past; draw $\xi_l$ and choose it in the event of
Construction~\ref{s7:consRound}(f); build $Q_l$, fix $\Dec_l$, lift it, and record
$\LentJV_l$ (step (4) of Construction~\ref{s6:consOrder}).
(iv) At rounds $2$ and $1$, run steps (1)--(2) of Construction~\ref{s6:consOrder};
here $J_2=J_1=\emptyset$.
\deps{\ref{s7:defSchedule}, \ref{s7:consRound}, \ref{s6:consOrder}, \ref{s4:lemTPV}, \ref{s4:lemPV}, \ref{s4:thmVXp}, \ref{s5:lemChild}, \ref{s5:lemDemoted}, \ref{s7:lemEXprime}, \ref{s7:lemPay}, \ref{s7:lemUHsplit}, \ref{s1:citMarkov}, \ref{s1:condG3}}
\fixes{\RTb{J5}, \RTb{I4}}
\end{lemma}

\begin{proof}
\emph{Step (i).} $\Xp$ is a non-negative random variable on a finite probability
space. By Markov's inequality (Cited result~\ref{s1:citMarkov}),
$\Prob(\Xp\le3\Ex\Xp)\ge2/3>0$, so the event is non-empty. Fix an outcome in it;
this gives (1). Stage~2 is a deterministic function of stage~1
(Definition~\ref{s7:defSchedule}).

\emph{Step (ii).} With stage~1 fixed, the good event of each part is determined
by stage-1 and stage-2 data together with that part's own stage-3 labels.
\begin{itemize}
\item For a standalone pre-part $Z$, Lemma~\ref{s4:lemTPV} is applied with $O_Z$
and $\vP=\Ret_Z$, both fixed at stage~2, and its good event is determined by
$(Z^0,O_Z,\Ret_Z)$ and the labels.
\item For a non-demoted light part, the good event of Lemma~\ref{s4:lemPV}
depends only on the own classes, $\Rt$, $\extph$ and the labels
(Lemma~\ref{s5:lemChild}).
\item For a demoted light part, the good event of Theorem~\ref{s4:thmVXp}
depends only on $O=X_Y$ and the labels (Lemma~\ref{s5:lemDemoted}).
\end{itemize}
Each of these events has probability at least $1/2-o(1)$, and by
condition~\ref{s1:condG3} every such $o(1)$ is at most $1/100$. So each event
is non-empty, whatever the stage-1 outcome; only this non-emptiness is used. The events of distinct parts involve disjoint families of labels,
so choosing each part's labels inside its own event gives one stage-3 outcome in
all of the events at once; no union bound is needed. None of these events
depends on the edge sets that the devices later receive: the conclusions of
Lemma~\ref{s4:lemTPV}, Lemma~\ref{s4:lemPV} and Theorem~\ref{s4:thmVXp} hold on
their good events for \emph{every} admissible edge set (every $E\supseteq E(O)$,
every $H_0$ containing the own classes, every graph containing $O$). So the
events do not depend on junk, on $\Erem$, on $H_0$ or on later rounds. This
gives (2).

\emph{Step (iii).} At round $l$ the past $\Past_l$ is fixed, and steps (1)--(3)
of Construction~\ref{s6:consOrder} are deterministic given it; in particular
$J_l$ is determined. The draw $\xi_l$ is independent of the past, and
$\copies_l$ and $\payrd_l$ are non-negative functions of $\xi_l$ with the past
fixed. By Markov's inequality (Cited result~\ref{s1:citMarkov}), the event of
Construction~\ref{s7:consRound}(f) has probability at least $1/2$. This holds for
every past and needs nothing beyond Markov's inequality. So a choice of $\xi_l$
in the event exists, whatever happened before. On this event, the explicit
bounds of (3) follow from Lemma~\ref{s7:lemUHsplit}(ii) and
Lemma~\ref{s7:lemPay}(c), which hold for \emph{every} past; this is the only
place where the uniformity in the past is used.

\emph{Step (iv).} Rounds $1$ and $2$ have no classed ports, since every port
there is fresh. So JS-LC produces no J-set, and there is no round step.
\end{proof}

\begin{proposition}[Cost on the chosen outcome]\label{s7:propCost}
On an outcome as in Lemma~\ref{s7:lemOneOutcome}, with the round step of
Construction~\ref{s7:consRound} as the J-consumer, the decomposition of $E(G)$
given by Theorem~\ref{s6:thmMIXC} has at most
\[
\Bigl(\frac{\Dstar}{2}+1085+\epsOne(\Dstar)\Bigr)n+2\sum_{l=3}^{R}\f(Q_l)
\]
objects, where
\[
\epsOne(\Dstar)\defeq\epsK(\Dstar)+169\,\epsA+\frac{252}{\Dstar}+1.5\,\epsCONC(\Dstar)+3\,\epsX(\Dstar)+\frac{9}{\Dstar},
\]
with $\epsA=31\eps/(\Cp\log\log\Dstar)$ as in Lemma~\ref{s2:lemTower}(e). The
terms are itemized in Table~\ref{s7:tabCost}. The constant is $745+338+2=1085$, and
$\epsOne$ is counted once. Consequently
$\Czero=\Dstar/2+1085+\epsOne(\Dstar)\le\Dstar/2+1091$ under \Gc{4}. No term
involves $\cEG$. The only Markov multipliers are the fixed $3$ (stage 1) and $4$
(per round). There is no term for VX-parts ($\mathcal{V}=\emptyset$).
\deps{\ref{s6:thmMIXC}, \ref{s6:thmCONCL}, \ref{s7:lemLift}, \ref{s7:lemPay}, \ref{s7:lemEXprime}, \ref{s7:lemOneOutcome}, \ref{s5:lemKRED}, \ref{s2:propParentless}, \ref{s2:propOV}, \ref{s7:consRound}, \ref{s2:lemTower}, \ref{s1:condG4}, \ref{s6:defLending}, \ref{s7:defXprime}}
\fixes{\RTb{J6}}
\end{proposition}

\begin{table}[ht]
\centering
\small
\begin{tabular}{|p{5.0cm}|p{6.2cm}|p{3.6cm}|}
\hline
\textbf{Item} & \textbf{Bound on the number of objects} & \textbf{Source}\\
\hline
Long cycles, $E_0$ and light parts (K-RED) & $(\Dstar/2+745+\epsK)n+[80\,\mathrm{dem}+369\,\mathrm{lp}]$ & Lemma~\ref{s5:lemKRED}\\
TPV in standalone pre-parts & $169(2+\epsA)n+[169\sum_l|\Lost_l|]$ & Thm.~\ref{s6:thmMIXC}(c); Props.~\ref{s2:propParentless}, \ref{s2:propOV}\\
JS-LC cycles & $252\,n/\Dstar+1.5\,\epsCONC(\Dstar)\,n$ & Thms.~\ref{s6:thmMIXC}(c), \ref{s6:thmCONCL}\\
$\Jlost$ items & $[\sum_l(M_l-1)|\Lost_l|]$ & Lemma~\ref{s7:lemPay}(a1)\\
Pooled and JV-bad items & $[\Xpool]$ & Lemma~\ref{s7:lemPay}(a2),(a3)\\
Unpaired fresh legs & $2n$ & Lemma~\ref{s7:lemPay}(b)\\
SDR failures and loops & $9n/\Dstar$ & Lemma~\ref{s7:lemPay}(d)\\
Lift of the quotients & $2\sum_l|\Dec_l|=2\sum_l\f(Q_l)$ & Lemma~\ref{s7:lemLift}\\
\hline
All bracketed terms together & $\le\XU+\Xpool\le\Xp\le3\,\Ex\Xp\le3\,\epsX(\Dstar)\,n$ & Lemmas~\ref{s7:lemEXprime}, \ref{s7:lemOneOutcome}\\
\hline
\end{tabular}
\caption{The cost line on the chosen outcome (Proposition~\ref{s7:propCost}).
Summing the rows gives $(\Dstar/2+745+338+2)n+\epsOne(\Dstar)n+2\sum_l\f(Q_l)$.}
\label{s7:tabCost}
\end{table}

\begin{proof}
By Lemma~\ref{s7:lemLift}(iv), the round step is a J-consumer. The outcome of
Lemma~\ref{s7:lemOneOutcome} has stage-3 labels in all good events. So
Theorem~\ref{s6:thmMIXC} applies to the run, $\delta$, the chosen stage-1 and
stage-3 outcomes and this J-consumer, and it yields a decomposition of $E(G)$. By
Theorem~\ref{s6:thmMIXC}(c), its number of objects is at most
\[
\Bigl(\frac{\Dstar}2+745+338\Bigr)n+\eps_M(\Dstar)\,n
+\Bigl[80\,\mathrm{dem}+369\,\mathrm{lp}+169\sum_l|\Lost_l|\Bigr]
+1.5\sum_{(Y,l)}m_{Y,l}+\sum_l|\Obj_l| ,
\]
with $\eps_M(\Dstar)=\epsK(\Dstar)+169\epsA+252/\Dstar$. This itemization
combines the K-RED bound of Lemma~\ref{s5:lemKRED}; the TPV bound
\[
169\sum_{l,Z}\bigl(|A_Z|+|F_Z|+|\Lost_Z|\bigr),
\]
with $\sum|F_Z|\le2n$ (Proposition~\ref{s2:propParentless}(i)) and
$\sum|A_Z|\le\epsA n$ (Proposition~\ref{s2:propOV}(K2) and
Lemma~\ref{s2:lemTower}(e)); and the JS-LC bound. By
Theorem~\ref{s6:thmCONCL}(iv),
$1.5\sum_{(Y,l)}m_{Y,l}\le1.5\,\epsCONC(\Dstar)\,n$.

The objects of $\Obj_l$ are the paid single edges and the lifted objects. By
Lemma~\ref{s7:lemPay}, the paid single edges number at most:
$(M_l-1)|\Lost_l|$ in (a1); the round-$l$ summand of $\Xpool$ in (a2) and (a3);
at most $2n$ unpaired fresh legs over all rounds; and $\payrd_l$ SDR-failure and
loop edges, with $\sum_l\payrd_l\le9n/\Dstar$ by Lemma~\ref{s7:lemPay}(d), since
the outcome satisfies Lemma~\ref{s7:lemOneOutcome}(3). By
Lemma~\ref{s7:lemLift}(iv), the lifted objects number at most $2\f(Q_l)$. Hence
\[
\sum_l|\Obj_l|\le\sum_l(M_l-1)|\Lost_l|+\Xpool+2n+\frac{9n}{\Dstar}+2\sum_l\f(Q_l).
\]
The bracketed terms and the terms $\sum_l(M_l-1)|\Lost_l|$ and $\Xpool$ together
equal
$80\,\mathrm{dem}+369\,\mathrm{lp}+\sum_l(169+M_l-1)|\Lost_l|+\Xpool\le\XU+\Xpool\le\Xp$.
By Lemma~\ref{s7:lemOneOutcome}(1) and Lemma~\ref{s7:lemEXprime},
$\Xp\le3\Ex\Xp\le3\epsX(\Dstar)n$. Adding everything gives at most
\[
\Bigl(\frac{\Dstar}2+745+338+2\Bigr)n
+\Bigl(\epsK+169\epsA+\frac{252}{\Dstar}+1.5\epsCONC+3\epsX+\frac9{\Dstar}\Bigr)n
+2\sum_l\f(Q_l),
\]
which is the claim, since $745+338+2=1085$. Every term of $\epsOne$ is a function
of $\Dstar$ alone, and $\epsOne(\Dstar)\le6$ is part
of \Gc{4}, which gives $\Czero\le\Dstar/2+1091$.
\end{proof}

\begin{theorem}[Theorem JV$^{+*}$]\label{s7:thmJVps}
Assume that $\Dstar$ satisfies \Gc{1}--\Gc{4}. Let $G$ be a graph with
$n\ge\Nzero$ vertices and $d_1\ge\Dstar$. Take any valid $\HBtp$ run on $G$
(Proposition~\ref{s2:propExists}), any designation $\delta$, and no VX-parts
($\mathcal{V}=\emptyset$). Then there are simple graphs $Q_3,\dots,Q_R$ without
isolated vertices such that
\begin{enumerate}
\item[(1)] $\f(G)\le\Czero n+2\sum_{l=3}^{R}\f(Q_l)$, where
$\Czero\defeq\Dstar/2+1085+\epsOne(\Dstar)$;
\item[(2)] $\sum_{l=3}^{R}|V(Q_l)|\le\thetaQ\,n$, where
$\thetaQ=\epsTwo(\Dstar)\le1/4$.
\end{enumerate}
Here $\epsOne$ and $\epsTwo$ are the explicit functions of $\Dstar$ defined in
Proposition~\ref{s7:propCost} and Lemma~\ref{s7:lemUHsplit}. Items at ultra hubs
cost no collision payments. They pay only the standard costs (pooled vertices,
JV-bad ports, SDR failures), and they add quotient vertices only through
Lemma~\ref{s7:lemUHsplit}.
\deps{\ref{s7:propCost}, \ref{s7:lemUHsplit}, \ref{s7:lemSimple}, \ref{s7:lemLift}, \ref{s7:lemOneOutcome}, \ref{s6:thmMIXC}, \ref{s2:propExists}, \ref{s7:lemPay}, \ref{s7:lemUltra}, \ref{s1:condGamma}}
\fixes{\RTa{M8}}
\end{theorem}

\begin{proof}
Take the outcome of Lemma~\ref{s7:lemOneOutcome} and let $Q_3,\dots,Q_R$ be the
quotients built on it. By Lemma~\ref{s7:lemSimple} each $Q_l$ is simple and has
no isolated vertices. By Lemma~\ref{s7:lemLift}(iv), the round step is a
J-consumer, so Theorem~\ref{s6:thmMIXC} gives a decomposition of $E(G)$. By
Proposition~\ref{s7:propCost} it has at most $\Czero n+2\sum_l\f(Q_l)$ objects,
which proves (1), since $\f(G)$ is at most the number of objects of any
decomposition. The outcome satisfies the hypotheses of
Lemma~\ref{s7:lemUHsplit}(iii): $\Xp\le3\Ex\Xp$, and the event of (f) holds at
every round. So Lemma~\ref{s7:lemUHsplit}(iv) gives (2). The statement about
ultra hubs is Lemma~\ref{s7:lemPay}(e) together with Lemma~\ref{s7:lemUltra}.
\end{proof}

%% ----------------------------------------------------------------------
\subsection{The layered quotient induction and the proof of the Main Theorem}

\begin{theorem}[Theorem HI$''$ (layered quotient induction)]\label{s7:thmHI}
Let $C\ge\Dstar/2$ and $\vartheta\in[0,1/2)$ be constants with the
following property. Every graph $G$ without isolated vertices, with $n\ge\Nzero$
vertices and $d_1=2|E(G)|/n\ge\Dstar$, admits finitely many simple graphs
$Q_1,\dots,Q_k$ ($k\ge0$) such that
\[
\f(G)\le C\,n+2\sum_{i=1}^{k}\f(Q_i)\qquad\text{and}\qquad\sum_{i=1}^{k}|V(Q_i)|\le\vartheta\,n .
\]
Then $\f(G)\le c\,|V(G)|$ for every graph $G$, where
$c\defeq\max\bigl(C/(1-2\vartheta),\,\Nzero/2\bigr)$.
\deps{\ref{s1:factAdd}, \ref{s1:defObject}, \ref{s1:defConstants}}
\fixes{\RTa{M6}}
\end{theorem}

\begin{proof}
Suppose the conclusion fails, and let $G$ be a counterexample with
$n\defeq|V(G)|$ minimal, so that $\f(G)>cn$. The graph without vertices has
$\f=0$, so $n\ge1$.

(1) \emph{$G$ has no isolated vertex.} If $v$ were isolated, then
$\f(G-v)=\f(G)$ by Fact~\ref{s1:factAdd}(d), so
$\f(G-v)>cn\ge c(n-1)=c|V(G-v)|$, and $G-v$ would be a smaller counterexample.

(2) \emph{$n\ge\Nzero$.} Otherwise $n-1<\Nzero$, and by Fact~\ref{s1:factAdd}(a),
$\f(G)\le|E(G)|\le n(n-1)/2<(\Nzero/2)\,n\le cn$.

(3) \emph{$d_1\ge\Dstar$.} Otherwise, by Fact~\ref{s1:factAdd}(a),
$\f(G)\le|E(G)|=d_1n/2<(\Dstar/2)\,n\le Cn\le cn$, because
$c\ge C/(1-2\vartheta)\ge C$. This is the trivial bound $\f(G)\le|E(G)|$; no
other estimate is used for graphs of small average degree.

(4) By (1)--(3), the hypothesis applies to $G$ and gives simple graphs
$Q_1,\dots,Q_k$ with $\f(G)\le Cn+2\sum_i\f(Q_i)$ and
$\sum_i|V(Q_i)|\le\vartheta n<n$. Each $Q_i$ has fewer vertices than $G$, so by
the minimality of $n$, $\f(Q_i)\le c|V(Q_i)|$.

(5) Hence $\f(G)\le Cn+2c\sum_i|V(Q_i)|\le Cn+2c\vartheta n
\le c(1-2\vartheta)n+2c\vartheta n=cn$, where we used $C\le c(1-2\vartheta)$.
This contradicts $\f(G)>cn$.

The constants $C$ and $\vartheta$ do not depend on $c$, and minimality is applied
only to the graphs $Q_i$, which have strictly fewer vertices than $G$.
\end{proof}

\begin{lemma}[The galactic conditions are satisfiable]\label{s7:lemGammaSat}
Let $\Nzero$ be as in Definition~\ref{s1:defConstants}(ii). Then:
\begin{enumerate}
\item[(i)] each item (a)--(f) of \Gc{1} (Condition~\ref{s1:condG1}) holds for
all sufficiently large real $\mu$; hence there is $\mu_1\ge1$ such that all of them
hold for every $\mu\ge\mu_1$, and \Gc{1} holds for every $\Dstar>2$ with
$\log_2\log_2\Dstar\ge\mu_1$;
\item[(ii)] $\epsOne(\Dstar)\to0$ and $\epsTwo(\Dstar)\to0$ as $\Dstar\to\infty$;
\item[(iii)] there is $D_0$ such that every $\Dstar\ge D_0$ satisfies
\Gc{1}--\Gc{4} simultaneously.
\end{enumerate}
In particular a constant $\Dstar$ as in Definition~\ref{s1:defConstants}(iii)
exists. Only this existence is used: neither $\mu_1$ nor $D_0$ is computed.
\deps{\ref{s1:condGamma}, \ref{s7:propCost}, \ref{s7:lemUHsplit}, \ref{s5:lemExpect}, \ref{s5:lemKRED}, \ref{s2:lemTower}, \ref{s7:lemEXprime}, \ref{s6:thmCONCL}, \ref{s1:defConstants}, \ref{s3:lemCOL}, \ref{s5:lemE1}, \ref{s3:lemCOLJVev}}
\end{lemma}

\begin{proof}
(i) We use inequalities of type (E) as defined in Lemma~\ref{s3:lemCOLJVev},
which has no hypothesis: by part~(i) of that lemma, every inequality of type (E)
holds for all sufficiently large $\mu$. For $\mu\ge1$ we have
$0\le\log_2\mu\le\log_2(\Aexp\mu)$, and:
item~(a) holds for every $\mu\ge2^8$;
item~(b), that is $\mu-14-\log_2\Aexp-3\log_2\mu\ge0$, follows from the
inequality $\mu-(14+\log_2\Aexp)-3\log_2(\Aexp\mu)\ge0$ of type (E);
item~(c) is the inequality $1.6\mu-2\Aexp\log_2(\Aexp\mu)\ge0$ of type (E);
item~(d), that is $\mu-8-2\log_2\mu\ge0$, follows from the inequality
$\mu-8-2\log_2(\Aexp\mu)\ge0$ of type (E);
item~(e), which is \Gstar, is (after taking $\log_2$ of both sides) the
inequality $36\mu-240-46\Aexp\log_2(\Aexp\mu)\ge0$ of type (E);
and every inequality of item~(f) holds for all sufficiently large $\mu$ by
Lemma~\ref{s3:lemCOLJVev}(iii).
So each of the finitely many inequalities of \Gc{1} holds for all $\mu$ at least
some threshold; let $\mu_1\ge1$ be at least all of these thresholds. If
$\Dstar>2$ and $\log_2\log_2\Dstar\ge\mu_1$, every $\mu\ge\log_2\log_2\Dstar$
satisfies $\mu\ge\mu_1$, so \Gc{1} holds.

(ii) By Proposition~\ref{s7:propCost} and Lemmas~\ref{s7:lemEXprime}
and~\ref{s7:lemUHsplit}, $\epsOne$ and $\epsTwo$ are finite sums of non-negative
explicit functions of $\Dstar$ with fixed coefficients:
$1/\Dstar$; $(6\log_2\Dstar+12)/\Dstar$;
$\epsA=31\eps/(\Cp\log\log\Dstar)$ (Lemma~\ref{s2:lemTower}(e));
$\epsCONC(\Dstar)$; $\epsU(\Dstar)$; $\epsK(\Dstar)$; and, in $\epsTwo$ only,
$F(\log_2\Dstar)$ with $F(x)=(3184+30400\log_2x)\,x^{-90.2}$ as in
Lemma~\ref{s7:lemUHsplit}(iv), which tends to $0$ because $F(x)\to0$ as
$x\to\infty$. The function $\epsCONC$ is the explicit expression
of Theorem~\ref{s6:thmCONCL}(iv). It is a combination of
$\logs\Dstar/\log\Dstar$, $\logs\Dstar/\log\log\Dstar$ and
$(2\logs\Dstar+2)/(\log\Dstar)^{1/2}$, each of which tends to $0$. The functions
$\epsU$ (Lemma~\ref{s5:lemExpect}) and $\epsK$ (Lemma~\ref{s5:lemKRED}) are, by
their defining statements, explicit functions of $\Dstar$ that tend to $0$ as
$\Dstar\to\infty$. Every other listed function tends to $0$. Hence
$\epsOne(\Dstar)\to0$ and $\epsTwo(\Dstar)\to0$.

(iii) Put $D_1\defeq2^{2^{\mu_1}}$; then $D_1\ge4>2$, as $\mu_1\ge1$. By (i),
\Gc{1} holds for every $\Dstar\ge D_1$, since then $\log_2\log_2\Dstar\ge\mu_1$.
\Gc{2}(a) holds for every $\Dstar\ge2^{117}$, and \Gc{2}(b),(c) are implied by
\Gc{1} (Condition~\ref{s1:condGamma}). \Gc{3} holds for every
$\Dstar\ge2^{(2\Nzero)^{1/\Cp}}$; here $\Nzero$ is fixed before $\Dstar$. By (ii)
there is $D''$ with $\epsOne(\Dstar)\le6$ and $\epsTwo(\Dstar)\le1/4$ for all
$\Dstar\ge D''$, which is \Gc{4}. Take
$D_0\defeq\max\bigl(D_1,2^{117},2^{(2\Nzero)^{1/\Cp}},D''\bigr)$.
\end{proof}

\begin{corollary}[Corollary JV$^{+*}$-EG; proof of the Main Theorem]\label{s7:thmMainProof}
Theorem~\ref{s1:thmMain} holds.
\deps{\ref{s7:thmHI}, \ref{s7:thmJVps}, \ref{s2:propExists}, \ref{s7:lemGammaSat}, \ref{s6:defDesign}, \ref{s1:condGamma}, \ref{s1:defConstants}}
\end{corollary}

\begin{proof}
Let $\Dstar$ be as in Theorem~\ref{s1:thmMain}, that is, a constant satisfying
\Gc{1}--\Gc{4}. Such constants exist by Lemma~\ref{s7:lemGammaSat}, and the
argument below applies to each of them. Let
$\Czero=\Dstar/2+1085+\epsOne(\Dstar)$, $\thetaQ=\epsTwo(\Dstar)$ and
$\cEG=\max\bigl(\Czero/(1-2\thetaQ),\Nzero/2\bigr)$, as in
Theorem~\ref{s1:thmMain}. We apply Theorem~\ref{s7:thmHI} with $C\defeq\Czero$
and $\vartheta\defeq\thetaQ$. Its assumptions on the constants hold:
$\Czero\ge\Dstar/2$, since $\epsOne\ge0$; and $0\le\thetaQ\le1/4<1/2$.

Let $G$ be a graph without isolated vertices, with $n\ge\Nzero$ vertices and
$d_1\ge\Dstar$. By Proposition~\ref{s2:propExists}, $G$ has a valid $\HBtp$
run. A designation exists: for every standalone pre-part $Z$ of a round $l\ge3$
and every $u\in Q_Z$, the set $\anc_l(u)$ is non-empty by the definition of
classed ports, and $Y(u)$ can be chosen in it by a fixed rule
(Definition~\ref{s6:defDesign}). With no VX-parts, Theorem~\ref{s7:thmJVps} gives
simple graphs $Q_3,\dots,Q_R$ with $\f(G)\le\Czero n+2\sum_l\f(Q_l)$ and
$\sum_l|V(Q_l)|\le\thetaQ n$. This is the hypothesis of
Theorem~\ref{s7:thmHI}, which therefore gives $\f(G)\le\cEG|V(G)|$ for every
graph $G$. In particular $\f(n)\le\cEG n$ for every $n$, so $\f(n)=O(n)$, where
$\cEG$ depends only on $\sigma,\Cp,\Aexp,\eps,\Nzero$ and $\Dstar$.

This completes the candidate proof; its status is described in Remark~\ref{s1:remStatus}.
\end{proof}

\begin{remark}[Why the argument is not circular]\label{s7:remNonCirc}
(1) $\Czero$ and $\thetaQ$ are fixed functions of $\Dstar$
(Proposition~\ref{s7:propCost}, Lemma~\ref{s7:lemUHsplit}), and $\Dstar$ is fixed
before any graph is considered. The constant $\cEG$ enters only as the
constant $c$ of Theorem~\ref{s7:thmHI} (throughout its statement and
proof); it plays no other role in any construction or bound of
Sections~\ref{s2:sec}--\ref{s7:sec}.
(2) Minimality is applied only to the quotients $Q_l$. They are simple and have
at most $\thetaQ n\le n/4<n$ vertices in total, and the induction uses about
them only the inequality $\f(Q_l)\le\cEG|V(Q_l)|$.
(3) The quotient $Q_l$ depends on the decompositions $\Dec_{l'}$ ($l'>l$), through
$J_l$ and through junk. This is harmless. Each $\Dec_{l'}$ is fixed at round $l'$
by a deterministic rule. The Markov choice at round $l$ succeeds with
probability at least $1/2$ whatever the past. Every bound of
Lemmas~\ref{s7:lemCC}, \ref{s7:lemUltra} and~\ref{s7:lemPay} holds for every
past.
(4) No other link calls itself. Lemma~TPV, Lemma~PV and Theorem~VX$^+$
finish their runs with the elementary long-cycle bound of
Fact~\ref{s1:factEG0}, which is proved directly and does not use the statement
being proved; all three are unconditional.
\deps{\ref{s7:thmHI}, \ref{s7:thmJVps}, \ref{s7:propCost}, \ref{s7:lemUHsplit}, \ref{s7:lemCC}, \ref{s7:lemUltra}, \ref{s7:lemPay}, \ref{s7:lemSimple}, \ref{s7:consRound}, \ref{s4:lemTPV}, \ref{s4:lemPV}, \ref{s4:thmVXp}, \ref{s1:condG4}, \ref{s1:factEG0}}
\end{remark}

%% ----------------------------------------------------------------------
\subsection*{What is not established}
\namedlabel{s7:ssecNotEstablished}{What is not established}

This section completes a \candidate{} proof. We record what has not been
established.

\begin{enumerate}
\item \emph{AI review only; no human verification.} No human expert has
verified any part of this manuscript. All checking so far was done by AI agents.
As a single text it has had six AI review rounds. Clean-room review~1
(eight lenses, referees reading only this text and \cite{BM}) found one major
gap, PV-ARC-END: the per-vertex arc-end bound of Lemma~\ref{s4:lemPV}(c) missed
arcs created by stripping in the Finish. The gap was repaired (fix
\textsf{CR1-PV}). Clean-room review~2 (two repair lenses, two lenses auditing
every site where a per-vertex multiplicity feeds a path-connectivity parameter
$t$, and a judge) found the repair correct and found no error at any such site;
its remaining findings were presentational. Clean-room reviews~3--5 each
returned verdict~(A): no fatal or major error, and presentational findings
only. A sixth round, a cross-model clean-room review (six section referees
running on a different AI model) together with three conceptual lenses
(proves-too-much, traces on concrete graph families, and key idea versus
barriers), also returned verdict~(A), with minor and cosmetic findings only,
all applied here. In blinded calibrations of this review process on copies of the text with
planted errors, the AI referees flagged every planted error that lay inside the
sections they were asked to review ($9$ of $9$ distinct plants), missed the one
plant placed outside those sections, and raised no false fatal or major alarm
in $21$ passes over the unmodified text.
Such AI review cannot exclude conceptual gaps that leave no local
inconsistency in the text, or blind spots shared by the model family that wrote
the proof and the one that reviewed it; verification by human experts is
required. Each link of
the chain was checked by two independent AI referees. Two AI red teams then
audited the chain and found no fatal or major error. Red team~1 listed the minor
fixes \RTa{M1}--\RTa{M12}; red team~2 listed \RTb{J1}--\RTb{J8} for the
composition and \RTb{I1}--\RTb{I16} for the inputs. Two further AI audits of the
routing engine of JS-LC also found no fatal or major error. All listed fixes are
applied in this text. Correlated blind spots between AI referees are possible.
The parts written or changed for the present version~6 of the text postdate
all of these reviews. They are described in Remark~\ref{s1:remStatus}(iii) and,
for item R6, in Section~\ref{s1:ssecDAG}. They have had one clean-room AI review
(gate G0, 2026-09-26: eight AI referee lenses and an AI judge), which covered
every changed statement and returned verdict~(A), no fatal or major error; its
minor and cosmetic fixes are applied here and have not been re-reviewed, and
the proof of the stated form of Haxell's theorem (Cited
result~\ref{s1:citHaxell}), added after G0, has had no referee yet. In
Table~\ref{s1:tabDAG}, the
statements that were added for version~6, or whose statement or proof was
rewritten for it, are marked ``one clean-room AI review (G0, 2026-09-26)'';
changes that concern only
single proof steps, dependency entries, pointers, names, titles or explanatory
text are not necessarily marked. G0, like every earlier review, was done by AI
agents; no human expert has refereed any part of the text, and the text remains
a \candidate{} proof. Version~6.1 (26 September 2026) applies the manuscript
patches found by the triage of the formalization in Lean (P2 triage): $M_l$ is
an integer, the fixed rules of Construction~\ref{s7:consRound} have declared
arguments, and several wording and proof-step repairs; no constant changes and
no statement used downstream is weakened (Lemma~\ref{s6:lemJSLC} is restated in
the form its only consumer uses; Remark~\ref{s1:remStatus}(iii)). These patches have had
no referee yet; the changed statements are marked ``changed in v6.1 (P2
triage): no referee yet'' in Table~\ref{s1:tabDAG}.
\item \emph{The weakest links} (as in Remark~\ref{s1:remStatus}(iv)).
(a) The routing engine inside JS-LC, that is, HCC-P, PAR, MED and EQ-LPT as used
for cherry systems at multiplicity $2M_l-2$ (Section~\ref{s6:sec}); it remains
listed first despite the two AI audits mentioned above.
(b) Lemma~14$^\tau$ together with CONC-L(iv), which modify the core Lemma~14 of
\cite{BM} (Sections~\ref{s2:sec} and~\ref{s6:sec}).
(c) The JV$^{+*}$ composition of this section, which is the newest link. Within
it, the most delicate points are: property \Jp{1} aggregated over all MIX-parts,
on which Lemma~MULT rests; the conditioning in Lemma~CC and in the ultra-hub
bound; the lift, including the cross-round use of junction edges through the
$\Lent$ sets; and the existence of one joint outcome.
(d) The statements changed by fix \textsf{CR1-PV} (Lemma~\ref{s4:lemPV}(c),
Lemma~\ref{s5:lemChild}(c), Lemma~\ref{s5:lemParent} Steps~5--6, the
parent-bad definition, Lemma~\ref{s3:lemCOL}(c) and rows~4 and~7 of
Table~\ref{s3:tabCOLJV}); these have had one AI re-review (clean-room
review~2).
(e) The factor $2q-1$ of Haxell's theorem (Cited result~\ref{s1:citHaxell}),
until it has been checked against the text of \cite{Hax95}; its only
application has slack of a factor of about $h/2$ (Lemma~\ref{s3:lemL9rho}),
and the proof of the stated form written out after the citation has had no
referee yet.
\item \emph{Refutation targets} (as in Remark~\ref{s1:remStatus}(v)). Any of the
following would refute the chain:
a valid $\HBtp$ run with $\sum_{(Y,l)}m_{Y,l}=\omega(n)$;
a valid run on which JS-LC emits two $\Jhub$ edges of one class at one hub in one
round (this would break Lemma~MULT);
a lifted sub-layer cycle that repeats a vertex;
an edge covered twice, in particular a junction edge used both by a lifted cycle
and by its owner's device;
any per-round cost term of order $\Theta(n)$ that is not a once-per-vertex
charge.
\item \emph{Status of the conjecture.} The \EG{} conjecture must not be recorded
as proved on the basis of this text. The argument here is a \candidate{} proof,
checked only by AI referees and not yet verified by human experts. It should be
regarded as a proof only after human experts have verified it.
\item \emph{No end-to-end computational check.} Only a toy end-to-end
bookkeeping simulation with stubbed analytic inputs exists; it proves nothing.
The other toy computations are for components: checkers of the
lift of the round step, a skeleton of the routing in JS-LC, and simulators of
Lemma~14$^\tau$. They test small instances in regimes far below the galactic
thresholds (Condition~\ref{s1:condGamma} makes $\Dstar$ of tower type), and they prove nothing.
\end{enumerate}
